\chapter{Levantamento e Análise de Requisitos}

\section{Estratégia e Métodos}

O levantamento de requisitos é uma das fases mais importantes, porque qualquer erro aqui pode causar problemas em todo o desenvolvimento. Para este projeto, usámos uma abordagem baseada nas boas práticas de engenharia de software (Sommerville, 2015), adaptada para a realidade de um clube local.

A nossa estratégia baseou-se em três pontos:
\begin{enumerate}
    \item \textbf{Usar várias fontes}: Não falámos apenas com uma pessoa nem olhámos apenas para um papel. Cruzámos informação de documentos, entrevistas e questionários para ter a certeza de que nada ficava esquecido.
    \item \textbf{Falar com quem usa o sistema}: Ouvimos as pessoas que vão realmente usar a aplicação no dia a dia, desde a gerência até aos próprios atletas.
    \item \textbf{Revisão constante}: Sempre que identificámos um requisito novo, voltámos a falar com o clube para confirmar se era mesmo aquilo que pretendiam.
\end{enumerate}

O primeiro passo foi identificar todos os grupos de pessoas que interagem com o clube, para percebermos quem seriam os nossos utilizadores principais (como atletas e treinadores) e quem seria afetado indiretamente (como a Autoridade Tributária por causa das faturas).

\begin{itemize}
    \item \textbf{\textit{Stakeholders} primários} (utilizadores diretos do sistema):
    \begin{itemize}
        \item Atletas -- praticantes das diversas modalidades do clube
        \item Treinadores -- responsáveis técnicos das equipas
        \item Encarregados de Educação -- responsáveis legais pelos atletas menores
        \item Secretaria -- funcionários administrativos do clube
        \item Gerência -- órgãos de direção e administração
    \end{itemize}

    \item \textbf{\textit{Stakeholders} secundários} (afetados indiretamente):
    \begin{itemize}
        \item Sócios não-praticantes -- membros associados sem atividade desportiva ativa
        \item Fornecedores e parceiros -- entidades externas com relações comerciais
        \item Autoridades desportivas -- federações e associações reguladoras
        \item Autoridade Tributária -- entidade reguladora fiscal (para efeitos de conformidade SAF-T)
    \end{itemize}
\end{itemize}

Para cada grupo de \textit{stakeholders}, foram identificadas as suas necessidades específicas, expectativas em relação ao sistema, e potenciais resistências à mudança. Esta análise prévia orientou a seleção dos métodos de levantamento mais adequados a cada perfil.

\caixaprompt{Gemini}{No contexto do CAP, quero que identifiques todos os stakeholders do clube.}{Para o Clube Amigos de Polvoreira (CAP), a identificação correta dos \textit{stakeholders} é fundamental para perceber quem irá interagir com o novo sistema e quem será afetado por ele. Tendo em conta a estrutura típica de um clube desportivo português, dividi os intervenientes em duas categorias principais:

\textbf{\textit{Stakeholders} Primários (Utilizadores Diretos):}
\begin{itemize}
    \item \textbf{Atletas:} Os praticantes das diversas modalidades do clube.
    \item \textbf{Treinadores:} Os responsáveis técnicos das equipas, que farão convocatórias e registo de lesões.
    \item \textbf{Encarregados de Educação:} Os responsáveis legais pelos atletas menores, fundamentais para a gestão de pagamentos e autorizações.
    \item \textbf{Secretaria:} Os funcionários administrativos que gerem inscrições, validação de atestados e processos diários.
    \item \textbf{Gerência:} Os órgãos de direção e administração que necessitam de uma visão global financeira e desportiva.
\end{itemize}

\textbf{\textit{Stakeholders} Secundários (Afetados Indiretamente):}
\begin{itemize}
    \item \textbf{Sócios não-praticantes:} Membros associados sem atividade desportiva ativa.
    \item \textbf{Fornecedores e parceiros:} Entidades externas com relações comerciais com o clube.
    \item \textbf{Autoridades desportivas:} Federações e associações reguladoras das competições.
    \item \textbf{Autoridade Tributária (AT):} Entidade reguladora fiscal, essencial para efeitos de conformidade e exportação do SAF-T.
\end{itemize}}



\subsection{Métodos de Levantamento Utilizados}

A equipa de desenvolvimento selecionou três métodos complementares de levantamento de requisitos, cada um com objetivos específicos e contribuições distintas para a compreensão global das necessidades do sistema:

\rowcolors{2}{rowblue}{white}
\begin{table}[H]
\centering
\begin{tabular}{|p{3.5cm}|p{5cm}|p{6cm}|}
\hline
\rowcolor{headerblue}\textcolor{white}{\textbf{Método}} & \textcolor{white}{\textbf{Objetivo Principal}} & \textcolor{white}{\textbf{Informação Obtida}} \\ \hline
Análise Documental & Compreender os processos atuais e os dados manipulados & Entidades de dados, fluxos de informação, regras de negócio implícitas \\ \hline
Entrevistas Direcionadas & Captar necessidades específicas e problemas vivenciados & Requisitos funcionais prioritários, dores dos utilizadores, expectativas qualitativas \\ \hline
Questionário Online & Validar requisitos em escala e identificar prioridades & Validação estatística, priorização de funcionalidades, identificação de consensos \\ \hline
\end{tabular}
\caption{Métodos de levantamento de requisitos e suas contribuições}
\end{table}
\rowcolors{2}{white}{white}

A complementaridade destes métodos foi essencial para garantir uma cobertura abrangente das necessidades do sistema. Enquanto a análise documental permitiu compreender ``o que existe'', as entrevistas revelaram ``o que falta'' e ``o que causa problemas'', e os questionários validaram ``o que é prioritário'' para a maioria dos utilizadores.

\caixaprompt{Gemini}{No contexto do CAP quero que me digas quais são os melhores métodos para o levantamento de requisitos para o nosso novo sistema de software.}{Para um ecossistema tão rico e diversificado como o do Clube Amigos de Polvoreira, recomendo uma abordagem mista (triangulação de métodos), inspirada nas práticas de Sommerville. Aqui estão os métodos mais indicados para o vosso caso:

\begin{itemize}
    \item \textbf{Análise Documental:} Fundamental para perceber "o que existe". Devem recolher os atestados médicos, relatórios de jogo em papel e os livros de recibos de quotas atuais para perceber exatamente que dados precisam de ser transpostos para o digital.
    
    \item \textbf{Entrevistas Direcionadas (Semiestruturadas):} Essencial para captar dores específicas. Recomendo entrevistar representantes chave: um elemento da direção (para a vertente financeira/SAF-T), um treinador (sobre convocatórias e lesões) e um atleta.
    
    \item \textbf{Questionários Online:} Sendo um clube com muitos associados, a melhor forma de validar problemas à escala (ex: dificuldades de comunicação em grupos de WhatsApp ou métodos de pagamento preferidos) é enviando um formulário via email para todos os atletas e encarregados de educação.
    
    \item \textbf{Observação Direta (Etnografia):} Passar uma tarde na secretaria do clube num dia de jogo ou de inscrições para observar os estrangulamentos do processo manual.
    
    \item \textbf{Focus Groups (Grupos de Foco):} Juntar numa sala treinadores de diferentes modalidades para debater como deve funcionar o sistema de reserva do pavilhão sem gerar conflitos.
    
    \item \textbf{Prototipagem Rápida:} Criar ecrãs falsos (\textit{mockups}) do sistema e mostrar à gerência para garantir que a usabilidade se adequa à literacia digital dos funcionários.
\end{itemize}

Recomendo que escolham 2 a 3 destes métodos para garantir uma boa cobertura sem esgotar o tempo útil do projeto.}

\subsection{Planeamento do Levantamento}

O processo de levantamento de requisitos foi organizado em três fases sequenciais, com duração total prevista de três semanas:

\begin{enumerate}
    \item \textbf{Fase 1 -- Análise Documental} (1 semana): Recolha e análise dos documentos físicos utilizados pelo clube, identificação de entidades de dados e mapeamento dos processos atuais.

    \item \textbf{Fase 2 -- Entrevistas} (1 semana): Realização de entrevistas semiestruturadas com representantes de cada grupo de \textit{stakeholders} primários, focando problemas atuais e funcionalidades desejadas.

    \item \textbf{Fase 3 -- Questionário e Consolidação} (1 semana): Distribuição do questionário online, análise dos resultados e consolidação de todos os requisitos identificados num documento de especificação preliminar.
\end{enumerate}

Na ata da primeira reunião de planeamento, ficou registada esta abordagem metodológica e a divisão de responsabilidades entre os membros da equipa:

\begin{figure}[H]
    \centering
    \makebox[\textwidth]{\includegraphics[width=0.9\textwidth]{ata1.drawio.png}}
    \caption{Ata reunião 1}
    \label{fig:ar-01}
\end{figure}

\subsection{Análise Documental}

A análise documental constituiu o ponto de partida do processo de levantamento de requisitos. Este método, frequentemente subestimado, revelou-se fundamental para compreender os processos operacionais existentes, identificar as entidades de dados que o sistema terá de gerir, e descobrir regras de negócio implícitas que dificilmente emergiriam apenas através de entrevistas.

\subsubsection{Objetivo e Abordagem}

O objetivo desta fase foi mapear o fluxo de informação atual do clube, identificando:
\begin{itemize}
    \item Os tipos de documentos processados regularmente;
    \item Os campos de dados registados em cada documento;
    \item As relações entre diferentes documentos e processos;
    \item As regras de validação e conformidade aplicáveis;
    \item Os pontos de falha e ineficiência do sistema atual.
\end{itemize}

A gerência do clube foi solicitada a disponibilizar exemplos representativos de cada categoria de documento utilizado na operação quotidiana. Após uma triagem inicial, foram selecionados três documentos-chave que cobrem as três áreas funcionais principais do clube: a vertente clínica, a vertente desportiva e a vertente financeira.

\subsubsection{Documentos Analisados}

\paragraph{Atestado Médico Desportivo}

A análise deste documento foi essencial para compreender os requisitos do módulo clínico do sistema. O atestado médico constitui um documento legalmente obrigatório para a prática desportiva federada, e a sua gestão adequada é crítica para a operação do clube.

Da análise do documento, foram extraídos os seguintes requisitos:
\begin{itemize}
    \item Necessidade de registar a data de emissão e a data de validade do atestado;
    \item Obrigatoriedade de associar cada atestado a um atleta específico;
    \item Requisito de armazenamento do documento digitalizado;
    \item Necessidade de um sistema de alertas automáticos para caducidade (30 dias antes, 7 dias antes, e no dia de expiração);
    \item Bloqueio automático de atletas com atestado caducado para convocatórias e participação em jogos.
\end{itemize}

\caixaprompt{Gemini}{Quero que, no contexto do clube, geres um exemplo de um atestado médico em HTML e de seguida me apresentes a lista de requisitos retirados desse documento.}{Com certeza! O ficheiro HTML foi gerado e estruturado com os campos padrão exigidos por lei para a prática desportiva federada. \textit{(Nota da equipa: O código e a pré-visualização do ficheiro HTML gerado encontram-se na secção de Anexos no final deste relatório).}}



Com base em cruzamento entre o HTML sintético, as notas da secretaria e os requisitos já listados acima, sistematizámos os seguintes pontos adicionais sob forma de RF:
\begin{itemize}
    \item \textbf{RF - Upload de Documento:} O sistema deve permitir aos atletas ou encarregados de educação efetuarem o \textit{upload} do atestado médico em formato digital (PDF ou imagem).
    \item \textbf{RF - Controlo de Validade:} O sistema deve extrair e registar a data de emissão do atestado médico, calculando a sua data de caducidade (habitualmente 1 ano após a emissão).
    \item \textbf{RF - Validação da Secretaria:} O sistema deve permitir aos funcionários da secretaria aprovar ou rejeitar o atestado submetido, validando a assinatura e a vinheta do médico.
    \item \textbf{RF - Alertas de Expiração:} O sistema deve notificar automaticamente o atleta e/ou encarregado de educação 30 dias antes de o atestado médico perder a validade.
    \item \textbf{RF - Bloqueio de Convocatória:} O sistema não deve permitir que os treinadores adicionem a convocatórias oficiais os atletas cujo atestado médico se encontre caducado ou reprovado.
\end{itemize}

\paragraph{Relatório Oficial de Jogo}

O relatório de jogo impresso permitiu mapear a vertente desportiva e estatística do sistema. Este documento é preenchido manualmente pela equipa técnica após cada partida e contém informação essencial para o acompanhamento do desempenho dos atletas.

Da análise do documento, foram identificadas as seguintes necessidades:
\begin{itemize}
    \item Registo dos atletas titulares e suplentes por jogo;
    \item Acompanhamento dos minutos jogados por cada atleta;
    \item Registo de eventos de jogo: golos marcados, assistências, cartões amarelos e vermelhos;
    \item Resultado final e dados do adversário;
    \item Observações técnicas do treinador.
\end{itemize}

\caixaprompt{Gemini}{Agora gera, no mesmo formato, um documento de relatório de jogo e mais uma vez extrai os requisitos do mesmo.}{Com certeza! O protótipo do Relatório Oficial de Jogo foi gerado em HTML, contendo os campos essenciais para o registo desportivo pós-partida. \textit{(Nota da equipa: À semelhança do documento anterior, o código e a pré-visualização encontram-se na secção de Anexos deste relatório).}}

A partir da estrutura e dos dados necessários para preencher este relatório de jogo, definiram-se como adenda os seguintes RF para o módulo desportivo:
\begin{itemize}
    \item \textbf{RF - Preenchimento do Relatório de Jogo:} O sistema deve permitir que o treinador (ou delegado) preencha o relatório após o apito final, registando o resultado final, a equipa adversária e a data/local da partida.
    \item \textbf{RF - Associação à Convocatória:} O sistema deve associar automaticamente o relatório de jogo à convocatória previamente gerada para essa partida, carregando a lista de atletas presentes.
    \item \textbf{RF - Registo de Estatísticas e Disciplina:} O sistema deve permitir o registo de eventos individuais por atleta durante o jogo, tais como golos marcados, assistências e cartões (amarelos/vermelhos).
    \item \textbf{RF - Registo de Lesões em Jogo:} O sistema deve possuir uma secção para assinalar lesões contraídas durante a partida, gerando um alerta automático para a secretaria e para o encarregado de educação (no caso de menores).
    \item \textbf{RF - Submissão e Arquivo:} Após a submissão do relatório pelo treinador, o documento deve ficar trancado a novas edições e disponível para consulta global pela gerência e secretaria.
\end{itemize}

Estas métricas definiram os requisitos para o módulo de estatísticas, permitindo a construção de perfis de desempenho individual e coletivo que a gerência poderá utilizar para apresentar a potenciais patrocinadores.

\paragraph{Fatura/Recibo de Quota}

A análise do documento de pagamento de quotas revelou os requisitos do módulo financeiro. O cruzamento do formato deste recibo com a lógica de negócio do clube permitiu compreender o ciclo de vida de um pagamento.

Da análise do documento, foram identificadas as seguintes necessidades:
\begin{itemize}
    \item Associação de identificadores únicos a cada pagamento;
    \item Registo de datas limite de pagamento por período (mensal, trimestral, anual);
    \item Gestão de estados de pagamento (pendente, pago, em atraso, reembolsado);
    \item Geração automática de faturas em conformidade com requisitos fiscais;
    \item Exportação para formato SAF-T conforme exigido pela Autoridade Tributária.
\end{itemize}

\caixaprompt{Gemini}{Agora gera, no mesmo formato, um documento de uma fatura de quota e mais uma vez extrai os requisitos do mesmo.}{Com certeza! O protótipo da Fatura/Recibo de Quota foi gerado em HTML, respeitando os campos essenciais de faturação de um clube desportivo. \textit{(Nota da equipa: O código e a pré-visualização deste ficheiro encontram-se na secção de Anexos do relatório).}}



Tendo consolidado estas restrições, avançámos para a seguinte versão sintética dos RF financeiros prioritários:

\begin{itemize}
    \item \textbf{RF - Processamento de Quotas:} O sistema deve permitir à secretaria registar o pagamento mensal ou anual de quotas por parte de sócios ou atletas.
    \item \textbf{RF - Emissão de Faturas e Recibos:} O sistema deve gerar automaticamente um documento em PDF (Fatura/Recibo) no momento da confirmação do pagamento, enviando-o para o email do utilizador.
    \item \textbf{RF - Integração SAF-T:} O sistema deve garantir que todos os recibos emitidos possuem os dados necessários (NIF, morada, taxa de IVA) para a correta exportação do ficheiro SAF-T para a Autoridade Tributária.
    \item \textbf{RF - Controlo de Dívidas:} O sistema deve identificar automaticamente atletas ou sócios com quotas em atraso, destacando-os no \textit{dashboard} da secretaria.
    \item \textbf{RF - Alertas de Pagamento:} O sistema deve enviar notificações automáticas (via email ou plataforma) aos encarregados de educação ou sócios nos 5 dias anteriores ao vencimento da quota.
    \item \textbf{RF - Restrição por Dívida:} O sistema deve avisar ou bloquear a inclusão de atletas em convocatórias caso tenham mais de 2 meses de quotas em atraso (conforme os regulamentos do clube).
\end{itemize}

\subsection{Entrevistas Direcionadas}

As entrevistas constituíram o método central de levantamento de requisitos, permitindo captar informação qualitativa que dificilmente emergiria através de outros métodos. Foram realizadas entrevistas semiestruturadas a três perfis essenciais do clube: um elemento da Gerência, um Treinador e um Atleta.

\subsubsection{Metodologia das Entrevistas}

Cada entrevista foi estruturada em três partes:

\begin{enumerate}
    \item \textbf{Contextualização} (5 minutos): Apresentação do projeto e dos objetivos da entrevista, esclarecimento sobre confidencialidade e obtenção de consentimento.

    \item \textbf{Questões estruturadas} (20 minutos): Conjunto de perguntas pré-definidas focadas em:
    \begin{itemize}
        \item Descrição das tarefas quotidianas relacionadas com o clube;
        \item Identificação de problemas e frustrações com os processos atuais;
        \item Funcionalidades desejadas para o novo sistema;
        \item Expectativas em termos de usabilidade e acessibilidade.
    \end{itemize}

    \item \textbf{Discussão aberta} (15 minutos): Espaço para o entrevistado levantar questões adicionais, propor funcionalidades não contempladas, e esclarecer dúvidas sobre o projeto.
\end{enumerate}

As entrevistas foram registadas em áudio (com autorização) e posteriormente transcritas para análise sistemática. De cada transcrição foram extraídos os requisitos implícitos e explícitos, categorizados por área funcional.

\subsubsection{Resultados por Perfil de \textit{Stakeholder}}

As conversas revelaram que a desorganização atual está a causar prejuízos financeiros, falhas de comunicação graves e até problemas de saúde física. Abaixo encontram-se sintetizados os principais pontos retirados de cada entrevista:

\subsubsection*{Perspetiva da Gerência}
\begin{itemize}
    \item \textbf{Perda de Investimento e Gestão Financeira:} O clube perdeu o patrocínio de uma empresa de bebidas energéticas simplesmente porque a gerência não conseguiu compilar e apresentar dados estruturados sobre o desempenho das equipas para justificar o investimento. Ficou clara a necessidade de \textit{dashboards} financeiros, gráficos de desempenho exportáveis e geração automática de ficheiros SAF-T.
    \item \textbf{Caos nas Quotas:} O modelo atual de pagamentos manuais resulta em esquecimentos frequentes por parte dos encarregados de educação. É urgente criar um perfil parental agregador com geração de referências Multibanco/MBWay e bloqueio automático de atletas com mensalidades em atraso.
    \item \textbf{Conflitos de Espaços:} Foi relatado um episódio em que as equipas de Futsal e Voleibol tentaram usar o pavilhão ao mesmo tempo devido a um erro num caderno de papel. O sistema tem de gerir o calendário de instalações.
\end{itemize}

\caixaprompt{Gemini}{Quero que geres uma entrevista feita por um membro da equipa de desenvolvimento a um Membro da gerência.}{Com certeza! Elaborei a transcrição de uma entrevista semiestruturada simulada entre um membro da equipa de desenvolvimento e um elemento da Direção (Gerência) do CAP. \textit{(Nota da equipa: A transcrição completa desta entrevista encontra-se na secção de Anexos no final deste relatório).}}



A partir das respostas consolidadas ao guião anterior, destacaram‑se sob forma de RF e notas funcionais os seguintes tópicos:

\begin{itemize}
    \item \textbf{RF - Dashboard de Gestão:} A direção necessita de um painel de controlo (\textit{dashboard}) que apresente estatísticas em tempo real, como o número de atletas inscritos por modalidade e o balanço financeiro mensal.
    \item \textbf{RF - Gestão de Perfis e Acessos:} O sistema deve suportar diferentes níveis de permissões. A gerência deve ter acesso administrativo global, garantindo que dados financeiros sensíveis não são visíveis para treinadores ou atletas.
    \item \textbf{RF - Conformidade Fiscal (SAF-T):} Foi sublinhado como requisito crítico e inegociável a capacidade de o sistema faturar e gerar o ficheiro SAF-T de forma automatizada para a Autoridade Tributária.
    \item \textbf{RF - Relatórios de Seguros e Atestados:} A gerência quer a capacidade de gerar um relatório rápido listando todos os atletas que se encontram com o seguro desportivo ou atestado médico caducado.
    \item \textbf{Dificuldade Atual Identificada:} O entrevistado destacou a perda de dezenas de horas mensais na verificação manual de transferências bancárias para reconciliação do pagamento de quotas.
\end{itemize}

\subsubsection*{Perspetiva da Equipa Técnica (Treinadores)}
\begin{itemize}
    \item \textbf{Falhas de Comunicação:} A dependência de grupos de WhatsApp para avisar sobre alterações de horários e convocatórias não funciona. Ocorreram situações em que a equipa foi para jogo com apenas um suplente porque os atletas não viram a mensagem. É obrigatório haver um sistema de convocatórias na plataforma com botões de confirmação (Sim/Não/Talvez).
    \item \textbf{Falta de Histórico Clínico:} Um atleta agravou uma lesão (rotura muscular) durante um treino porque o treinador não foi informado da queixa física que tinha ocorrido em casa. O sistema precisa de um módulo onde a equipa técnica possa consultar os atletas inaptos em tempo real.
    \item \textbf{Estatísticas e Material:} Há a necessidade de gráficos radar para comparar capacidades de atletas da mesma posição, além de um sistema rápido para registar assiduidade, ocorrências de jogo (golos, cartões) e requisitar material à secretaria.
\end{itemize}

\caixaprompt{Gemini}{Quero que geres uma entrevista feita por um membro da equipa de desenvolvimento a um Treinador do clube.}{Com certeza! Elaborei a transcrição de uma entrevista semiestruturada entre a equipa de desenvolvimento e um Treinador do CAP. \textit{(Nota da equipa: A transcrição completa encontra-se no ficheiro `EntrevistaTreinador.txt` na secção de Anexos deste relatório).}}
\subsection{Estratégia de Validação}
Sob esta ótica, sintetizaram-se os seguintes RF de apoio à equipa técnica:

\begin{itemize}
    \item \textbf{RF - Gestão de Convocatórias:} O treinador precisa de um módulo onde possa selecionar os atletas para o jogo do fim de semana, enviando uma notificação automática para os mesmos.
    \item \textbf{RF - Bloqueio de Atletas Irregulares:} O sistema deve impedir que o treinador convoque atletas que tenham o atestado médico caducado ou quotas em atraso, evitando penalizações desportivas.
    \item \textbf{RF - Preenchimento do Relatório de Jogo:} O treinador necessita de registar de forma célere os eventos da partida (golos, cartões, resultado final) diretamente no telemóvel ou \textit{tablet} após o jogo.
    \item \textbf{RF - Registo de Assiduidades e Lesões:} Capacidade de registar as presenças nos treinos semanais e reportar lesões imediatas ao departamento médico/secretaria.
    \item \textbf{Dificuldade Atual Identificada:} O treinador queixou-se de que perde muito tempo a tentar saber pela secretaria quem está "apto" a jogar devido a papelada atrasada, operando muitas vezes "às cegas".
\end{itemize}

\subsubsection*{Perspetiva dos Atletas}
\begin{itemize}
    \item \textbf{Atestados Médicos Caducados:} Um problema crítico identificado foi a falta de avisos sobre a validade dos exames médicos desportivos, o que levou a que um atleta ficasse de fora de uma competição importante. A plataforma tem de alertar automaticamente para a caducidade destes documentos.
    \item \textbf{Autonomia Clínica e Financeira:} Os atletas queremos poder registar queixas físicas autonomamente no sistema para que cheguem de imediato ao treinador. Querem também acesso fácil aos horários, histórico de pagamento de quotas e faturas automáticas.
    \item \textbf{Perfil Individual Desportivo:} Pediram uma área pessoal onde possam ver a evolução do seu peso, altura e estatísticas ao longo da época. Mostraram também interesse numa biblioteca multimédia para vídeos táticos partilhados pelo treinador e num espaço para o palmarés individual (prémios e distinções).
\end{itemize}

\caixaprompt{Gemini}{Quero que geres uma entrevista feita por um membro da equipa de desenvolvimento a um Atleta do clube.}{Sem problema! Criei a simulação de uma entrevista focada na experiência de um Atleta do CAP (e do seu Encarregado de Educação). \textit{(Nota da equipa: A transcrição na íntegra está disponível no ficheiro `EntrevistaAtleta.txt` nos Anexos do relatório).}}

Dos tópicos validados pelo atleta/EE, destacaram‑se estas necessidades funcionais nucleares:

\begin{itemize}
    \item \textbf{RF - Consulta de Calendário e Convocatórias:} O atleta precisa de aceder a uma área pessoal onde possa ver a data, hora e local dos próximos jogos, bem como confirmar se está convocado.
    \item \textbf{RF - Pagamento e Histórico de Quotas:} É essencial que o atleta consiga consultar o estado das suas mensalidades, obter as referências de pagamento (ex: MBWay, Multibanco) e fazer o \textit{download} das faturas/recibos anteriores.
    \item \textbf{RF - Submissão de Documentação:} O sistema deve permitir ao atleta enviar o atestado médico digitalizado e atualizar os seus dados pessoais (ex: alteração de morada ou contacto de emergência).
    \item \textbf{RF - Notificações Proativas:} O atleta deseja receber alertas no telemóvel ou por email quando uma quota está prestes a expirar ou quando o treinador publica uma convocatória de última hora.
    \item \textbf{Dificuldade Atual Identificada:} A maior frustração apontada foi a necessidade de ir presencialmente à secretaria do clube (ficar em filas) apenas para entregar o atestado médico em papel no início de cada época.
\end{itemize}

\subsection{Questionário Online}

\subsubsection{Objetivo e Metodologia}

Para complementar a análise documental e as entrevistas individuais, foi desenvolvido um questionário online com o objetivo de escalar a recolha de informação e obter dados quantitativos de uma amostra significativa de \textit{stakeholders}. Este método permitiu validar estatisticamente os problemas identificados nas fases anteriores e priorizar as funcionalidades a desenvolver com base em consensos da comunidade do clube.

O questionário foi construído utilizando a plataforma Microsoft Forms e distribuído eletronicamente a todos os membros do clube com endereço de correio eletrónico registado. O período de recolha de respostas foi de uma semana, tendo sido obtido um total de 47 respostas válidas.

\subsubsection{Estrutura do Questionário}

O formulário foi organizado em quatro secções temáticas:

\begin{enumerate}
    \item \textbf{Identificação do perfil}: Tipo de utilizador (atleta, encarregado de educação, treinador, funcionário), escalão etário, modalidade(s) praticada(s).

    \item \textbf{Comunicação e organização atual}: Avaliação da eficácia dos canais de comunicação existentes (grupos de WhatsApp, chamadas telefónicas, comunicação presencial), frequência de falhas de informação, impacto dessas falhas na operação do clube.

    \item \textbf{Gestão financeira}: Facilidade no pagamento de quotas, frequência de atrasos involuntários, preferência por métodos de pagamento, necessidade de acesso a histórico de pagamentos.

    \item \textbf{Funcionalidades desejadas}: Lista de funcionalidades potenciais para priorização (sistema de convocatórias, gestão de atestados médicos, reserva de instalações, estatísticas desportivas, etc.), espaço para sugestões adicionais em resposta aberta.
\end{enumerate}

\caixaprompt{Gemini}{Quero que geres a estrutura de um questionário Google Forms para ser enviado a todos os stakeholders, e depois vamos gerar as perguntas 1 prompt de cada vez.}{Excelente abordagem. Para garantir que o questionário capte as necessidades de todos os intervenientes do Clube Amigos de Polvoreira de forma organizada, recomendo dividir o Google Forms nas seguintes secções principais, utilizando a lógica de ramificação (\textit{branching}) do formulário:

\begin{itemize}
    \item \textbf{Secção 1: Perfil do Inquirido:} Pergunta inicial de seleção múltipla (Atleta, Encarregado de Educação, Treinador, Sócio ou Secretaria/Gerência) para direcionar o utilizador.
    \item \textbf{Secção 2: Avaliação Geral (Todos):} Questões sobre os canais de comunicação atuais (WhatsApp, email, presencial) e as maiores dores na obtenção de informação do clube.
    \item \textbf{Secção 3: Módulos Específicos (Ramificados):}
    \begin{itemize}
        \item \textit{Para Treinadores:} Dificuldades nas convocatórias, assiduidades e relatórios de jogo.
        \item \textit{Para Pais/Atletas:} Experiência com o pagamento de quotas, envio de atestados e acesso ao calendário de jogos.
    \end{itemize}
    \item \textbf{Secção 4: O Novo Sistema:} Pergunta de resposta aberta para sugestão de funcionalidades prioritárias na nova plataforma.
\end{itemize}}

\vspace{0.3cm}
\hrule
\vspace{0.2cm}
\textit{\textbf{Nota da Equipa de Desenvolvimento:} Por motivos de concisão do relatório, optámos por não documentar nesta secção a totalidade das iterações seguintes. Conforme planeado, o desenvolvimento das perguntas exatas e opções de resposta foi realizado de forma iterativa, através de um \textit{prompt} dedicado a cada pergunta individual. }



\subsubsection{Análise dos Resultados}

Os resultados obtidos confirmaram os padrões identificados nas entrevistas:

\begin{itemize}
    \item \textbf{87\%} dos inquiridos classificaram a comunicação atual como ``insuficiente'' ou ``muito insuficiente'';
    \item \textbf{72\%} reportaram pelo menos uma situação nos últimos três meses em que falharam um compromisso por falta de informação atempada;
    \item \textbf{91\%} manifestaram interesse em receber notificações automáticas sobre convocatórias e alterações de horário;
    \item \textbf{65\%} indicaram preferência por pagamento de quotas através de referência Multibanco ou MB Way;
    \item \textbf{78\%} dos encarregados de educação solicitaram acesso a um portal onde possam consultar a situação financeira e os documentos dos seus educandos.
\end{itemize}

O cruzamento destes dados quantitativos com os testemunhos das entrevistas permitiu estabelecer uma hierarquia clara de prioridades para o desenvolvimento do sistema, colocando os módulos de comunicação/notificações e gestão financeira no topo das funcionalidades a implementar.

\subsubsection{Exemplo de Questão}

Na figura seguinte é apresentado um exemplo de uma das questões que constavam no formulário enviado:

\begin{figure}[H]
    \centering
    \makebox[\textwidth]{\includegraphics[width=0.9\textwidth]{forms.png}}
    \caption{Exemplo de questão do questionário enviado aos \textit{stakeholders}}
    \label{fig:forms-exemplo}
\end{figure}

\subsection{Síntese e Validação dos Requisitos Levantados}

A triangulação das três fontes de informação (análise documental, entrevistas e questionário) permitiu construir uma visão abrangente e validada das necessidades do sistema. Os requisitos identificados foram compilados num documento de especificação preliminar, que foi apresentado à gerência do clube para validação antes de prosseguir para a fase de especificação formal.

Nesta reunião de validação, foram discutidos os seguintes pontos:
\begin{itemize}
    \item Confirmação da cobertura de todas as áreas funcionais identificadas;
    \item Priorização das funcionalidades a desenvolver nas primeiras iterações;
    \item Esclarecimento de ambiguidades e requisitos conflituantes;
    \item Identificação de restrições técnicas e orçamentais que possam afetar o âmbito do projeto.
\end{itemize}

A validação foi bem-sucedida, tendo a gerência aprovado o documento de requisitos e autorizado o avanço para a fase de especificação formal. As secções seguintes apresentam os requisitos funcionais e não funcionais resultantes deste processo de levantamento.

\section{Requisitos Funcionais}\label{sec:requisitos_funcionais}

\noindent\textit{Atalho: Saltar para \hyperref[sec:requisitos_nao_funcionais]{Requisitos Não-Funcionais} ou para o \hyperref[chap:especificacao_modelacao]{Próximo Capítulo (Modelação)}.}\par\medskip

% ---------------------------------------------------------
% REQUISITO 1
% ---------------------------------------------------------
O primeiro requisito funcional, designado por Gestão de Perfis e Utilizadores (RF-01), insere-se na área de Gestão Administrativa e assume-se como o pilar estrutural da plataforma. Este requisito surge da necessidade da gerência e da secretaria do clube poderem criar, editar e arquivar os perfis de todos os membros associados à instituição, atribuindo-lhes os respetivos níveis de permissão. Para que qualquer uma destas ações seja realizada, verifica-se como pré-condição que o utilizador responsável pela operação se encontre previamente autenticado no sistema com os privilégios adequados. 

A ordem de operações estipulada dita que, ao registar ou editar um perfil, o sistema deve primeiro validar a unicidade de dados críticos como o email e o NIF. Caso o sistema detete que estes identificadores já existem na base de dados, a operação de registo é interrompida e o erro é sinalizado. Em caso de sucesso, cumpre-se a pós-condição da operação: o novo perfil fica registado e apto a aceder ao sistema, sendo as credenciais iniciais enviadas automaticamente por correio eletrónico. 

Em relação à operação de arquivamento, caso um atleta ou membro do staff abandone o clube, o sistema reage revogando imediatamente as suas capacidades de acesso. No entanto, assegura-se que toda a informação desportiva e financeira ligada a esse perfil é mantida intacta. A relevância deste requisito para a aplicação é crítica, uma vez que é ele que permite identificar as pessoas que interagem com o sistema e que ações podem realizar, servindo de base para todos os outros módulos.

{\small
\rowcolors{2}{rowblue}{white}
\begin{xltabular}{\textwidth}{|l|X|}
\hline
\rowcolor{headerblue}\textcolor{white}{\textbf{Campo}} & \textcolor{white}{\textbf{Descrição}} \\ \hline
\endfirsthead

\multicolumn{2}{c}{{\tablename\ \thetable{} -- Continuação}} \\
\hline
\rowcolor{headerblue}\textcolor{white}{\textbf{Campo}} & \textcolor{white}{\textbf{Descrição}} \\ \hline
\endhead

\hline \multicolumn{2}{|r|}{\textit{Continua na próxima página...}} \\ \hline \endfoot

\hline
\caption{Especificação do Requisito Funcional: RF-01}
\label{tab:rf-01} \\
\endlastfoot

\textbf{ID} & RF-01 \\ \hline
\textbf{Nome do Requisito} & Gestão de Utilizadores e Perfis \\ \hline
\textbf{Área do Sistema} & Gestão Global de Perfis \\ \hline
\textbf{Ator Principal} & Atletas, Encarregados de Educação, Treinadores, Staff e Gerência \\ \hline
\textbf{Descrição} & O sistema deve permitir a criação, edição e consulta de perfis para todos os intervenientes do clube. A visualização e edição das informações (pessoais, desportivas e médicas) devem ser controladas de forma dinâmica, dependendo estritamente do nível de acesso e do cargo (perfil de segurança) do utilizador autenticado. \\ \hline
\textbf{Prioridade} & Alta \\ \hline
\textbf{Origem} & Entrevista --- Gerência do CAP (Reunião 1) \\ \hline
\textbf{Analista} & Tiago Martins \\ \hline
\textbf{Pré-condições} & 1. O sistema de controlo de acessos baseado em papéis (RBAC - \textit{Role-Based Access Control}) tem de estar configurado. \newline 2. Para a visualização ou edição de perfis de terceiros, o utilizador tem de estar autenticado e possuir os privilégios necessários (ex: um Treinador tem de estar associado à equipa do Atleta que pretende consultar). \\ \hline
\textbf{Pós-condições} & 1. Os dados do perfil são persistidos na base de dados com a devida encriptação em campos sensíveis (ex: passwords, dados médicos). \newline 2. É gerado um registo de auditoria (\textit{log}) interno documentando quem criou ou alterou o perfil e a respetiva data/hora. \newline 3. O perfil fica disponível ou restrito nas listagens globais do clube, conforme as regras de negócio associadas ao seu papel. \\ \hline
\textbf{Critérios de Aceitação} & 1. O sistema deve garantir que um atleta/encarregado apenas consegue editar os seus próprios dados pessoais. \newline 2. A interface deve bloquear a visualização de dados financeiros (ex: quotas) aos Treinadores, limitando o acesso a esses campos estritamente à Gerência/Secretaria e ao próprio Atleta. \newline 3. A criação de contas com perfis de nível superior (Gerência, Treinadores) deve ser obrigatoriamente aprovada ou realizada por um Administrador do sistema. \\
\end{xltabular}
}
\rowcolors{2}{white}{white}


% ---------------------------------------------------------
% REQUISITO 2
% ---------------------------------------------------------
O requisito que se segue, designado por Gestão de Equipas e Plantéis (RF-02), dá continuidade à estruturação do sistema. Uma vez assegurada a criação dos perfis, surge a necessidade prática de agrupar os atletas e a equipa técnica nas respetivas equipas e escalões competitivos. Este requisito determina que a direção desportiva e a secretaria devem ser capazes de criar, editar e inativar equipas, bem como de gerir a alocação dos atletas aos respetivos plantéis. Para realizar estas operações, verifica-se como pré-condição que o utilizador se encontre autenticado com os privilégios de gestão adequados e que os perfis dos atletas a alocar já existam na base de dados. Durante o processo de associação de um atleta a uma equipa, o sistema tem o dever de executar uma validação automática da data de nascimento do atleta face aos limites etários do escalão selecionado. Caso detete uma incompatibilidade, a operação é rejeitada e o utilizador é notificado do erro. Em caso de sucesso, cumpre-se a pós-condição: o plantel é atualizado e o atleta fica oficialmente apto a ser integrado nas atividades da equipa, o que se reflete de imediato nos módulos dependentes, como as convocatórias e a gestão de presenças em treinos. A relevância deste requisito é alta, dado que a divisão do capital humano em equipas representa o núcleo de toda a atividade desportiva da instituição.

{\small
\rowcolors{2}{rowblue}{white}
\begin{xltabular}{\textwidth}{|l|X|}
\hline
\rowcolor{headerblue}\textcolor{white}{\textbf{Campo}} & \textcolor{white}{\textbf{Descrição}} \\ \hline
\endfirsthead

\multicolumn{2}{c}{{\tablename\ \thetable{} -- Continuação}} \\
\hline
\rowcolor{headerblue}\textcolor{white}{\textbf{Campo}} & \textcolor{white}{\textbf{Descrição}} \\ \hline
\endhead

\hline \multicolumn{2}{|r|}{\textit{Continua na próxima página...}} \\ \hline \endfoot

\hline
\caption{Especificação do Requisito Funcional: RF-02}
\label{tab:rf-02} \\
\endlastfoot

\textbf{ID} & RF-02 \\ \hline
\textbf{Nome do Requisito} & Gestão de Perfis Parentais \\ \hline
\textbf{Área do Sistema} & Gestão de Utilizadores e Perfis \\ \hline
\textbf{Ator Principal} & Encarregado de Educação / Gerência \\ \hline
\textbf{Descrição} & O sistema deve permitir a criação de um perfil de "Encarregado de Educação" e a sua associação formal a um ou múltiplos perfis de atletas (tipicamente menores de idade), centralizando a gestão de autorizações, acompanhamento desportivo e responsabilidades financeiras. \\ \hline
\textbf{Prioridade} & Alta \\ \hline
\textbf{Origem} & Questionário Online (47 respostas) + Entrevista Atleta \\ \hline
\textbf{Analista} & Pedro Ferreira \\ \hline
\textbf{Pré-condições} & 1. O perfil do Encarregado de Educação e os perfis dos respetivos atletas já se encontram criados no sistema. \newline 2. O processo de vinculação exige validação por parte da Secretaria/Gerência ou a introdução de um código/token de segurança único fornecido pelo clube para garantir a legitimidade da associação familiar. \\ \hline
\textbf{Pós-condições} & 1. O Encarregado de Educação herda automaticamente as permissões de visualização e edição (quando aplicável) dos dados dos atletas associados. \newline 2. As contas correntes (quotas e pagamentos) dos vários educandos são consolidadadas sob a responsabilidade financeira deste perfil parental. \newline 3. O sistema redireciona os alertas e notificações (ex: convocatórias, lesões, quotas em atraso) dos menores para os contactos do Encarregado de Educação. \\ \hline
\textbf{Critérios de Aceitação} & 1. O sistema deve permitir que um encarregado de educação consulte a informação de vários educandos alternando entre eles de forma fluida na mesma interface. \newline 2. Deve ser possível gerar uma única referência de pagamento (MBWay/Multibanco) que englobe as quotas de todos os dependentes associados. \newline 3. Caso um atleta atinja a maioridade, o sistema deve permitir a desvinculação administrativa ou a transição de responsabilidades, caso seja solicitado. \\
\end{xltabular}
}
\rowcolors{2}{white}{white}


% ---------------------------------------------------------
% REQUISITO 3
% ---------------------------------------------------------
O requisito que se segue, designado por Gestão de Inscrições e Renovações (RF-03), dá continuidade ao processo de estruturação administrativa do clube. Uma vez que os perfis estão criados e os atletas alocados aos respetivos plantéis, é estritamente necessário garantir que a situação legal e desportiva de cada praticante está regularizada para a época em curso. Este requisito determina que os atletas, ou os seus encarregados de educação, podem submeter o pedido de inscrição ou renovação através da plataforma, anexando a documentação desportiva e clínica exigida. Tratando-se de um processo que envolve a submissão e verificação de ficheiros confidenciais, foi estipulado que a documentação submetida transita para um estado de análise e que a validação e aprovação final destas inscrições apenas pode ser levada a cabo pelos funcionários da secretaria do clube. A conclusão e o cumprimento deste requisito revelam-se de extrema importância para o sistema, uma vez que apenas os atletas com a inscrição devidamente aprovada poderão ser oficialmente inseridos nas convocatórias e participar em competições, servindo também de ponto de partida para a geração automática do plano de pagamento de quotas associado a essa época.

{\small
\rowcolors{2}{rowblue}{white}
\begin{xltabular}{\textwidth}{|l|X|}
\hline
\rowcolor{headerblue}\textcolor{white}{\textbf{Campo}} & \textcolor{white}{\textbf{Descrição}} \\ \hline
\endfirsthead

\multicolumn{2}{c}{{\tablename\ \thetable{} -- Continuação}} \\
\hline
\rowcolor{headerblue}\textcolor{white}{\textbf{Campo}} & \textcolor{white}{\textbf{Descrição}} \\ \hline
\endhead

\hline \multicolumn{2}{|r|}{\textit{Continua na próxima página...}} \\ \hline \endfoot

\hline
\caption{Especificação do Requisito Funcional: RF-03}
\label{tab:rf-03} \\
\endlastfoot

\textbf{ID} & RF-03 \\ \hline
\textbf{Nome do Requisito} & Sistema de Comunicação Interna (Chat) \\ \hline
\textbf{Área do Sistema} & Comunicação e Notificações \\ \hline
\textbf{Ator Principal} & Atletas, Encarregados de Educação, Treinadores e Gerência \\ \hline
\textbf{Descrição} & O sistema deve disponibilizar um módulo centralizado de mensagens seguras que permita a comunicação direta (1-para-1) ou em grupo entre os vários intervenientes do clube, respeitando a sua estrutura hierárquica e desportiva. \\ \hline
\textbf{Prioridade} & Alta \\ \hline
\textbf{Origem} & Entrevista --- Treinador do CAP (Reunião 2) \\ \hline
\textbf{Analista} & André Santos \\ \hline
\textbf{Pré-condições} & 1. O utilizador tem de estar autenticado. \newline 2. O sistema de turmas/equipas tem de estar configurado para restringir a lista de contactos (ex: um atleta só pode iniciar conversa com os seus treinadores diretos ou com a secretaria, evitando spam global). \\ \hline
\textbf{Pós-condições} & 1. A mensagem é armazenada de forma persistente e segura na base de dados, criando um histórico institucional e auditável. \newline 2. O destinatário (ou destinatários, no caso de grupos) recebe uma notificação imediata através da plataforma (e/ou push notification no telemóvel). \newline 3. O estado de leitura da mensagem é sincronizado entre os intervenientes. \\ \hline
\textbf{Critérios de Aceitação} & 1. A interface deve apresentar a lista de contactos organizados logicamente (ex: "A minha Equipa", "Equipa Técnica", "Secretaria"). \newline 2. O sistema deve permitir que um treinador ou elemento da gerência crie grupos de conversação (ex: "Convocados Jogo Sábado"). \newline 3. O módulo deve suportar o envio de pequenos anexos (ex: imagens ou PDFs para justificações de faltas ou partilha de táticas). \\
\end{xltabular}
}
\rowcolors{2}{white}{white}


% ---------------------------------------------------------
% REQUISITO 4
% ---------------------------------------------------------
O quarto requisito funcional, designado por Notificações e Alertas (RF-04), assume um papel transversal e de suporte a toda a dinâmica da plataforma. Este requisito surge da necessidade de garantir que o sistema gera e envia avisos automáticos aos diversos utilizadores (como atletas, encarregados de educação, treinadores e staff) sempre que ocorram eventos que exijam a sua atenção. Entre estes eventos incluem-se a publicação de novas convocatórias, alterações de última hora em horários ou locais de treino, e avisos de quotas ou pagamentos em atraso. Para que as notificações sejam emitidas com sucesso, verifica-se como pré-condição que o utilizador destinatário tenha os seus dados de contacto devidamente preenchidos no perfil e que o evento que despoleta o aviso tenha sido acionado no sistema. Assim que a ação ocorre, o sistema atua de forma proativa para cumprir a pós-condição: a informação chega imediatamente ao utilizador, seja através de alertas visuais internos na própria interface da aplicação ou através de comunicações externas, como o correio eletrónico. A relevância deste requisito para a aplicação é bastante elevada, dado que assegura uma comunicação fluida, contínua e atempada entre a estrutura do clube e os seus membros, mitigando falhas de informação no dia a dia.

{\small
\rowcolors{2}{rowblue}{white}
\begin{xltabular}{\textwidth}{|l|X|}
\hline
\rowcolor{headerblue}\textcolor{white}{\textbf{Campo}} & \textcolor{white}{\textbf{Descrição}} \\ \hline
\endfirsthead

\multicolumn{2}{c}{{\tablename\ \thetable{} -- Continuação}} \\
\hline
\rowcolor{headerblue}\textcolor{white}{\textbf{Campo}} & \textcolor{white}{\textbf{Descrição}} \\ \hline
\endhead

\hline \multicolumn{2}{|r|}{\textit{Continua na próxima página...}} \\ \hline \endfoot

\hline
\caption{Especificação do Requisito Funcional: RF-04}
\label{tab:rf-04} \\
\endlastfoot

\textbf{ID} & RF-04 \\ \hline
\textbf{Nome do Requisito} & Motor de Notificações Automáticas \\ \hline
\textbf{Área do Sistema} & Comunicação e Notificações \\ \hline
\textbf{Ator Principal} & Atletas, Encarregados de Educação, Treinadores e Gerência \\ \hline
\textbf{Descrição} & O sistema deve possuir um motor transversal capaz de gerar e enviar alertas automáticos (via push notification, email ou na própria plataforma) relativos a eventos críticos do clube, tais como alterações de horários de treinos, publicação de convocatórias, avisos de quotas em atraso e alertas de caducidade de exames médicos. \\ \hline
\textbf{Prioridade} & Alta \\ \hline
\textbf{Origem} & Questionário Online (47 respostas) \\ \hline
\textbf{Analista} & Pedro Ferreira \\ \hline
\textbf{Pré-condições} & 1. O utilizador tem de ter os seus dados de contacto (email/telefone) corretos no perfil e as permissões de notificação ativas na aplicação ou browser. \newline 2. Os gatilhos (eventos) têm de ocorrer no sistema (ex: um treinador altera a hora do treino, ou o servidor deteta que faltam 30 dias para um atestado caducar). \\ \hline
\textbf{Pós-condições} & 1. A notificação é processada e despachada para os canais adequados de cada utilizador afetado. \newline 2. O sistema gera um registo (log) de auditoria comprovando que o aviso foi enviado com sucesso (essencial para evitar desculpas como "não fui avisado que o exame caducava"). \\ \hline
\textbf{Critérios de Aceitação} & 1. Alterações de horários de treinos ou espaços físicos devem gerar uma notificação imediata (real-time) para a equipa correspondente. \newline 2. O sistema deve emitir alertas automáticos para a renovação do exame médico a 30 e a 15 dias da data de caducidade. \newline 3. O sistema deve enviar um lembrete automático e cordial aos encarregados de educação/atletas nos dias seguintes ao vencimento de uma quota não paga. \\
\end{xltabular}
}
\rowcolors{2}{white}{white}


% ---------------------------------------------------------
% REQUISITO 5
% ---------------------------------------------------------
O quinto requisito funcional, designado por Gestão de Treinos e Registo de Presenças (RF-05), introduz a vertente de Gestão Desportiva na plataforma. Este requisito surge da necessidade prática da equipa técnica organizar as sessões de trabalho semanais e monitorizar de forma rigorosa a assiduidade dos seus atletas. Para a realização destas operações, verifica-se como pré-condição que o utilizador se encontre autenticado com o perfil de Treinador ou Coordenador Técnico, e que as equipas e os respetivos plantéis já estejam previamente configurados e alocados no sistema. Durante a operação de registo, o treinador acede ao evento de treino agendado e assinala o estado de cada jogador (presente, falta justificada, falta injustificada ou lesionado). O sistema encarrega-se de processar estas informações e, após a submissão, cumpre-se a pós-condição: os dados são guardados de forma persistente e as estatísticas individuais de assiduidade no perfil de cada atleta são atualizadas em tempo real. A relevância deste requisito para a aplicação é fulcral, uma vez que o planeamento dos treinos e o acompanhamento contínuo da presença dos atletas representam a atividade central do clube no terreno. Esta funcionalidade substitui os tradicionais registos em papel e fornece dados estatísticos essenciais para apoiar as futuras decisões técnicas, servindo de base para o módulo de convocações para os jogos oficiais.

{\small
\rowcolors{2}{rowblue}{white}
\begin{xltabular}{\textwidth}{|l|X|}
\hline
\rowcolor{headerblue}\textcolor{white}{\textbf{Campo}} & \textcolor{white}{\textbf{Descrição}} \\ \hline
\endfirsthead

\multicolumn{2}{c}{{\tablename\ \thetable{} -- Continuação}} \\
\hline
\rowcolor{headerblue}\textcolor{white}{\textbf{Campo}} & \textcolor{white}{\textbf{Descrição}} \\ \hline
\endhead

\hline \multicolumn{2}{|r|}{\textit{Continua na próxima página...}} \\ \hline \endfoot

\hline
\caption{Especificação do Requisito Funcional: RF-05}
\label{tab:rf-05} \\
\endlastfoot

\textbf{ID} & RF-05 \\ \hline
\textbf{Nome do Requisito} & Gestão de Modalidades e Épocas \\ \hline
\textbf{Área do Sistema} & Gestão Desportiva e Operacional \\ \hline
\textbf{Ator Principal} & Gerência \\ \hline
\textbf{Descrição} & O sistema deve permitir a criação, edição e inativação de modalidades desportivas oferecidas pelo clube, bem como a gestão completa do ciclo das épocas desportivas (abertura e fecho), automatizando a transição e a promoção dos atletas para os respetivos escalões etários no início de cada nova época. \\ \hline
\textbf{Prioridade} & Alta \\ \hline
\textbf{Origem} & Entrevista --- Gerência do CAP (Reunião 1) \\ \hline
\textbf{Analista} & Tiago Martins \\ \hline
\textbf{Pré-condições} & 1. O utilizador encontra-se autenticado com privilégios de nível de Gerência ou de Administrador do Sistema. \newline 2. Não existem dependências financeiras ou desportivas ativas a bloquear o encerramento da época anterior (ex: todos os jogos da época atual devem estar dados como concluídos). \\ \hline
\textbf{Pós-condições} & 1. As novas modalidades e escalões ficam imediatamente disponíveis nos formulários de inscrição para atletas e encarregados de educação. \newline 2. No fecho de época, os registos de presenças, convocações e relatórios técnicos dessa época transitam para um estado de "apenas leitura" (histórico). \newline 3. O sistema recalcula as idades desportivas e aloca provisoriamente os atletas ativos aos escalões da nova época. \\ \hline
\textbf{Critérios de Aceitação} & 1. A inativação de uma modalidade não deve apagar o histórico de dados desportivos ou financeiros do clube, devendo apenas ocultá-la de novas inscrições. \newline 2. O processo de transição de época deve apresentar à gerência uma listagem de pré-visualização das subidas de escalão antes da sua efetivação final na base de dados. \newline 3. O sistema deve permitir a configuração de escalões etários personalizados por modalidade (ex: Sub-15 no Futebol pode ter regras de idade diferentes de Cadetes na Natação). \\
\end{xltabular}
}
\rowcolors{2}{white}{white}


% ---------------------------------------------------------
% REQUISITO 6
% ---------------------------------------------------------
O sexto requisito funcional, designado por Gestão de Convocatórias (RF-06), dá continuidade à vertente de Gestão Desportiva da plataforma. Este requisito surge da necessidade da equipa técnica definir e comunicar oficialmente quais os atletas selecionados para representar o clube num determinado jogo, seja ele de caráter oficial ou amigável. Para a realização desta operação, verifica-se como pré-condição que o utilizador se encontre autenticado com o perfil de Treinador ou Coordenador Técnico e que o evento correspondente ao jogo já esteja criado no calendário do sistema. Durante o processo de seleção, a interface deve auxiliar a decisão do treinador, cruzando os dados do plantel com as informações de assiduidade aos treinos e o histórico clínico (para evitar a convocação de atletas lesionados ou a cumprir castigo). Uma vez confirmada e submetida a seleção, cumpre-se a pós-condição: a lista de convocados é guardada na base de dados de forma persistente e o sistema aciona de imediato o envio de alertas automáticos aos atletas selecionados e aos respetivos encarregados de educação. A relevância deste requisito para a aplicação é bastante elevada, uma vez que representa o culminar do ciclo de trabalho semanal da equipa. A desmaterialização das convocatórias garante que a informação chega de forma rápida e inequívoca aos intervenientes envolvidos, minimizando falhas de comunicação e facilitando toda a organização logística para os dias de jogo.

{\small
\rowcolors{2}{rowblue}{white}
\begin{xltabular}{\textwidth}{|l|X|}
\hline
\rowcolor{headerblue}\textcolor{white}{\textbf{Campo}} & \textcolor{white}{\textbf{Descrição}} \\ \hline
\endfirsthead

\multicolumn{2}{c}{{\tablename\ \thetable{} -- Continuação}} \\
\hline
\rowcolor{headerblue}\textcolor{white}{\textbf{Campo}} & \textcolor{white}{\textbf{Descrição}} \\ \hline
\endhead

\hline \multicolumn{2}{|r|}{\textit{Continua na próxima página...}} \\ \hline \endfoot

\hline
\caption{Especificação do Requisito Funcional: RF-06}
\label{tab:rf-06} \\
\endlastfoot

\textbf{ID} & RF-06 \\ \hline
\textbf{Nome do Requisito} & Planeamento de Horários e Treinos \\ \hline
\textbf{Área do Sistema} & Gestão Desportiva e Operacional \\ \hline
\textbf{Ator Principal} & Treinadores / Gerência (Criação) e Atletas / Encarregados de Educação (Consulta) \\ \hline
\textbf{Descrição} & O sistema deve permitir aos treinadores ou membros da gerência criar, editar e cancelar sessões de treino. Qualquer marcação ou alteração efetuada deve refletir-se em tempo real nos calendários de consulta dos atletas afetados. \\ \hline
\textbf{Prioridade} & Alta \\ \hline
\textbf{Origem} & Entrevista --- Treinador do CAP (Reunião 2) \\ \hline
\textbf{Analista} & André Santos \\ \hline
\textbf{Pré-condições} & 1. O utilizador criador encontra-se autenticado com as permissões adequadas (Treinador associado à equipa ou Gerência). \newline 2. A época, a modalidade e a equipa/turma em questão encontram-se ativas. \newline 3. Os espaços físicos do clube (campos, pavilhões) já estão parametrizados na base de dados. \\ \hline
\textbf{Pós-condições} & 1. A sessão de treino é persistida na base de dados, reservando automaticamente o espaço físico selecionado para aquele horário. \newline 2. O calendário individual na área dos atletas correspondentes é imediatamente atualizado. \newline 3. Caso a alteração ocorra num prazo curto (ex: menos de 24h), o sistema aciona o motor de notificações para alertar os envolvidos sobre a mudança. \\ \hline
\textbf{Critérios de Aceitação} & 1. O formulário de marcação deve suportar a criação de eventos recorrentes (ex: treinos todas as terças e quintas às 19h00 durante toda a época). \newline 2. O sistema tem de bloquear ativamente a gravação se detetar uma sobreposição de horários no mesmo espaço físico com outra equipa. \newline 3. A interface de calendário dos atletas deve destacar visualmente (com cores ou ícones) os treinos que sofreram alterações recentes de horário ou local. \\
\end{xltabular}
}
\rowcolors{2}{white}{white}


% ---------------------------------------------------------
% REQUISITO RF-07
% ---------------------------------------------------------

O requisito RF-07, denominado Gestão de Presenças e Faltas, integra-se na área da Gestão Desportiva e Operacional e visa digitalizar o processo de controlo de assiduidade nas sessões de treino. Nas organizações desportivas tradicionais, este registo é frequentemente realizado em papel ou em ferramentas dispersas, impossibilitando a análise estatística da assiduidade ao longo da época e dificultando a deteção precoce de padrões de absentismo.

O sistema concretiza este requisito ao gerar automaticamente a lista de chamada com base nos atletas inscritos, permitindo ao treinador, através de uma interface otimizada para dispositivos móveis, marcar rapidamente a presença ou falta de cada elemento. O histórico resultante fica imediatamente disponível na área pessoal do atleta e nos painéis de gestão, garantindo que encarregados de educação podem acompanhar a assiduidade dos seus educandos.

{\small
\rowcolors{2}{rowblue}{white}
\begin{xltabular}{\textwidth}{|l|X|}
\hline
\rowcolor{headerblue}\textcolor{white}{\textbf{Campo}} & \textcolor{white}{\textbf{Descrição}} \\ \hline
\endfirsthead
\multicolumn{2}{c}{{\tablename\ \thetable{} -- Continuação}} \\
\hline
\rowcolor{headerblue}\textcolor{white}{\textbf{Campo}} & \textcolor{white}{\textbf{Descrição}} \\ \hline
\endhead
\hline \multicolumn{2}{|r|}{\textit{Continua na próxima página...}} \\ \hline \endfoot
\hline
\caption{Especificação do Requisito Funcional: RF-07}
\label{tab:rf-07} \\
\endlastfoot
\textbf{ID} & RF-07 \\ \hline
\textbf{Nome do Requisito} & Gestão de Presenças e Faltas \\ \hline
\textbf{Área do Sistema} & Gestão Desportiva e Operacional \\ \hline
\textbf{Ator Principal} & Treinadores e Atletas/EE \\ \hline
\textbf{Descrição} & O sistema deve permitir ao treinador realizar o registo digital da assiduidade de cada sessão de treino, marcando a presença ou falta de cada atleta convocado. As faltas devem ser categorizadas por tipologia (justificada ou injustificada) e o historial de assiduidade deve ser acessível tanto ao treinador como ao próprio atleta ou encarregado de educação. \\ \hline
\textbf{Prioridade} & Alta \\ \hline
\textbf{Origem} & Entrevista --- Treinador do CAP (Reunião 2) \\ \hline
\textbf{Analista} & André Santos \\ \hline
\textbf{Pré-condições} & 1. Autenticado com perfil de Treinador associado à equipa/turma. \newline 2. Sessão de treino criada e agendada no calendário. \newline 3. Atletas inscritos na modalidade e escalão corretos. \\ \hline
\textbf{Pós-condições} & 1. Registo de assiduidade persistido e sessão passa a ``Realizada/Concluída''. \newline 2. Estatísticas de assiduidade recalculadas automaticamente em \textit{background}. \newline 3. Histórico visível na área pessoal do atleta e painéis gerenciais. \\ \hline
\textbf{Critérios de Aceitação} & 1. Interface responsiva e otimizada para uso em telemóveis (lista com botões de alternância rápida). \newline 2. Sistema carrega automaticamente a lista de chamada. \newline 3. Atleta/encarregado visualiza motivo e tipologia de faltas. \\
\end{xltabular}
}
\rowcolors{2}{white}{white}

% ---------------------------------------------------------
% REQUISITO RF-08
% ---------------------------------------------------------

O requisito RF-08, designado Gestão de Convocatórias e Disponibilidade, insere-se na área da Gestão Desportiva e Operacional e tem como objetivo centralizar o processo de seleção e comunicação de convocatórias para eventos desportivos. A gestão eficaz de convocatórias é fundamental para garantir que os treinadores conseguem planear a participação das suas equipas com informação fiável sobre a disponibilidade de cada atleta.

Através deste módulo, o treinador pode selecionar atletas de uma lista filtrada, visualizando imediatamente quaisquer impedimentos (lesões, suspensões, indisponibilidades declaradas). Os atletas e encarregados de educação recebem uma notificação imediata e podem confirmar ou recusar a presença diretamente na plataforma, sendo obrigatória a justificação em caso de recusa. O treinador dispõe de um painel em tempo real com o estado de cada convocado.

{\small
\rowcolors{2}{rowblue}{white}
\begin{xltabular}{\textwidth}{|l|X|}
\hline
\rowcolor{headerblue}\textcolor{white}{\textbf{Campo}} & \textcolor{white}{\textbf{Descrição}} \\ \hline
\endfirsthead
\multicolumn{2}{c}{{\tablename\ \thetable{} -- Continuação}} \\
\hline
\rowcolor{headerblue}\textcolor{white}{\textbf{Campo}} & \textcolor{white}{\textbf{Descrição}} \\ \hline
\endhead
\hline \multicolumn{2}{|r|}{\textit{Continua na próxima página...}} \\ \hline \endfoot
\hline
\caption{Especificação do Requisito Funcional: RF-08}
\label{tab:rf-08} \\
\endlastfoot
\textbf{ID} & RF-08 \\ \hline
\textbf{Nome do Requisito} & Gestão de Convocatórias e Disponibilidade \\ \hline
\textbf{Área do Sistema} & Gestão Desportiva e Operacional \\ \hline
\textbf{Ator Principal} & Treinadores, Atletas e EE \\ \hline
\textbf{Descrição} & O sistema deve permitir ao treinador criar e divulgar listas de convocatória para eventos desportivos (jogos, torneios, concentrações), permitindo que cada atleta ou encarregado de educação confirme ou recuse a disponibilidade através da plataforma, com indicação obrigatória do motivo em caso de recusa. \\ \hline
\textbf{Prioridade} & Alta \\ \hline
\textbf{Origem} & Entrevista --- Treinador do CAP (Reunião 2) \\ \hline
\textbf{Analista} & André Santos \\ \hline
\textbf{Pré-condições} & 1. Autenticado com perfil de Treinador. \newline 2. Evento desportivo criado no calendário com data, hora e local. \newline 3. Atletas selecionáveis pertencem ao plantel da equipa. \\ \hline
\textbf{Pós-condições} & 1. Lista de convocados persistida na BD. \newline 2. Motor de notificações envia alerta imediato. \newline 3. Respostas de disponibilidade acessíveis em painel de controlo em tempo real. \\ \hline
\textbf{Critérios de Aceitação} & 1. Seleção de atletas por lista ou filtros, visualizando impedimentos. \newline 2. Interface com botões de ação claros e motivo obrigatório se negativo. \newline 3. \textit{Dashboard} com confirmados, ausentes e pendentes. \\
\end{xltabular}
}
\rowcolors{2}{white}{white}

% ---------------------------------------------------------
% REQUISITO RF-09
% ---------------------------------------------------------

O requisito RF-09, denominado Planos de Treino e Recursos Táticos, engloba-se na área da Gestão Desportiva e Operacional e responde à necessidade de organizar e partilhar conteúdos técnicos entre a equipa técnica e os atletas. A centralização de planos de treino, vídeos de análise e esquemas táticos numa plataforma única elimina a dependência de canais de comunicação informais e garante o controlo de acesso à informação sensível de cada equipa.

O treinador pode fazer o \textit{upload} de ficheiros multimédia (vídeos MP4, esquemas em PDF, imagens) e associá-los diretamente a sessões de treino ou épocas específicas. Os atletas da equipa acedem imediatamente ao conteúdo em modo de leitura, sendo o acesso estritamente isolado entre diferentes equipas para salvaguardar a confidencialidade tática.

{\small
\rowcolors{2}{rowblue}{white}
\begin{xltabular}{\textwidth}{|l|X|}
\hline
\rowcolor{headerblue}\textcolor{white}{\textbf{Campo}} & \textcolor{white}{\textbf{Descrição}} \\ \hline
\endfirsthead
\multicolumn{2}{c}{{\tablename\ \thetable{} -- Continuação}} \\
\hline
\rowcolor{headerblue}\textcolor{white}{\textbf{Campo}} & \textcolor{white}{\textbf{Descrição}} \\ \hline
\endhead
\hline \multicolumn{2}{|r|}{\textit{Continua na próxima página...}} \\ \hline \endfoot
\hline
\caption{Especificação do Requisito Funcional: RF-09}
\label{tab:rf-09} \\
\endlastfoot
\textbf{ID} & RF-09 \\ \hline
\textbf{Nome do Requisito} & Planos de Treino e Recursos Táticos \\ \hline
\textbf{Área do Sistema} & Gestão Desportiva e Operacional \\ \hline
\textbf{Ator Principal} & Treinadores e Atletas \\ \hline
\textbf{Descrição} & O sistema deve permitir ao treinador criar e partilhar planos de treino estruturados e recursos táticos multimédia (vídeos de análise, esquemas táticos em PDF, imagens) com os atletas da sua equipa, organizando esses conteúdos por sessão ou época desportiva. \\ \hline
\textbf{Prioridade} & Média \\ \hline
\textbf{Origem} & Entrevista --- Treinador do CAP (Reunião 2) \\ \hline
\textbf{Analista} & André Santos \\ \hline
\textbf{Pré-condições} & 1. Treinador autenticado e associado à equipa. \newline 2. Módulo de armazenamento configurado com espaço disponível. \\ \hline
\textbf{Pós-condições} & 1. Plano/ficheiro armazenado e associado à equipa ou sessão. \newline 2. Atletas com acesso imediato em modo leitura. \newline 3. Notificação opcional enviada aos atletas. \\ \hline
\textbf{Critérios de Aceitação} & 1. \textit{Upload} aceita MP4, PDF, imagens com limite de tamanho. \newline 2. Associação direta a sessão agendada. \newline 3. Isolamento de informação entre equipas. \\
\end{xltabular}
}
\rowcolors{2}{white}{white}

% ---------------------------------------------------------
% REQUISITO RF-10
% ---------------------------------------------------------

O requisito RF-10, denominado Registo de Ocorrências e Disciplina, insere-se na área da Gestão Desportiva e Operacional e permite o registo estruturado de eventos ocorridos em contexto de jogo e de incidentes disciplinares. A distinção clara entre estatísticas desportivas (golos, assistências, cartões) e ocorrências de cariz disciplinar é essencial para a gestão rigorosa do comportamento dos atletas e para o cumprimento das obrigações perante as federações.

O sistema alimenta automaticamente as estatísticas individuais de cada atleta com os eventos registados e aplica bloqueios preventivos em convocatórias sempre que uma suspensão disciplinar esteja em vigor. O acesso à informação é controlado por perfil: um atleta ou encarregado de educação apenas visualiza os registos respeitantes ao próprio, enquanto a equipa técnica e a gerência têm acesso à visão global do plantel.

{\small
\rowcolors{2}{rowblue}{white}
\begin{xltabular}{\textwidth}{|l|X|}
\hline
\rowcolor{headerblue}\textcolor{white}{\textbf{Campo}} & \textcolor{white}{\textbf{Descrição}} \\ \hline
\endfirsthead
\multicolumn{2}{c}{{\tablename\ \thetable{} -- Continuação}} \\
\hline
\rowcolor{headerblue}\textcolor{white}{\textbf{Campo}} & \textcolor{white}{\textbf{Descrição}} \\ \hline
\endhead
\hline \multicolumn{2}{|r|}{\textit{Continua na próxima página...}} \\ \hline \endfoot
\hline
\caption{Especificação do Requisito Funcional: RF-10}
\label{tab:rf-10} \\
\endlastfoot
\textbf{ID} & RF-10 \\ \hline
\textbf{Nome do Requisito} & Registo de Ocorrências e Disciplina \\ \hline
\textbf{Área do Sistema} & Gestão Desportiva e Operacional \\ \hline
\textbf{Ator Principal} & Treinadores/Gerência e Atletas/EE \\ \hline
\textbf{Descrição} & O sistema deve permitir o registo estruturado de ocorrências no contexto desportivo, distinguindo entre eventos de jogo (golos, assistências, cartões, substituições) e ocorrências disciplinares (comportamentos antidesportivos, sanções federativas). As ocorrências devem alimentar automaticamente as estatísticas individuais e podem gerar suspensões que bloqueiam o atleta em convocatórias futuras. \\ \hline
\textbf{Prioridade} & Média \\ \hline
\textbf{Origem} & Análise Documental (Relatório Oficial de Jogo) \\ \hline
\textbf{Analista} & Tiago Martins \\ \hline
\textbf{Pré-condições} & 1. Autenticado com privilégios necessários. \newline 2. Partida/competição criada no calendário. \newline 3. Atleta visado com perfil ativo. \\ \hline
\textbf{Pós-condições} & 1. Ocorrência guardada no histórico do atleta. \newline 2. Estatísticas de desempenho atualizadas. \newline 3. Suspensão disciplinar bloqueia inclusão em convocatórias. \\ \hline
\textbf{Critérios de Aceitação} & 1. Distinção clara entre ``Eventos de Jogo'' e ``Ocorrências Disciplinares''. \newline 2. Atleta/EE vê apenas as suas próprias ocorrências. \newline 3. Sanção grave permite anexar relatório PDF. \\
\end{xltabular}
}
\rowcolors{2}{white}{white}

% ---------------------------------------------------------
% REQUISITO RF-11
% ---------------------------------------------------------

O requisito RF-11, designado Gestão do Exame Médico Desportivo (EMD), insere-se na área da Saúde e Bem-Estar e representa um requisito de cumprimento obrigatório sob a perspetiva legal e federativa. O certificado de aptidão física, emitido por médico competente, é condição sine qua non para que um atleta possa participar em competições oficiais, pelo que a sua gestão centralizada e o controlo da validade são críticos para a operacionalidade do clube.

O sistema permite o carregamento seguro do documento e o acompanhamento do estado do atleta ao longo do ciclo de vida do certificado. São emitidos alertas proativos antes da data de expiração, e, caso o prazo de validade seja ultrapassado sem renovação, o sistema altera automaticamente o estado do atleta para ``Inapto'', impedindo a sua inclusão em convocatórias até que um novo certificado válido seja submetido e aprovado pela gerência.

{\small
\rowcolors{2}{rowblue}{white}
\begin{xltabular}{\textwidth}{|l|X|}
\hline
\rowcolor{headerblue}\textcolor{white}{\textbf{Campo}} & \textcolor{white}{\textbf{Descrição}} \\ \hline
\endfirsthead
\multicolumn{2}{c}{{\tablename\ \thetable{} -- Continuação}} \\
\hline
\rowcolor{headerblue}\textcolor{white}{\textbf{Campo}} & \textcolor{white}{\textbf{Descrição}} \\ \hline
\endhead
\hline \multicolumn{2}{|r|}{\textit{Continua na próxima página...}} \\ \hline \endfoot
\hline
\caption{Especificação do Requisito Funcional: RF-11}
\label{tab:rf-11} \\
\endlastfoot
\textbf{ID} & RF-11 \\ \hline
\textbf{Nome do Requisito} & Gestão do Exame Médico Desportivo \\ \hline
\textbf{Área do Sistema} & Saúde e Bem-Estar \\ \hline
\textbf{Ator Principal} & Atletas/EE e Gerência \\ \hline
\textbf{Descrição} & O sistema deve centralizar a gestão do Exame Médico Desportivo (EMD), permitindo o \textit{upload} seguro do certificado médico, o controlo da data de validade e a emissão de alertas proativos de renovação. O EMD é um requisito obrigatório federativo sem o qual o atleta não pode participar em competições oficiais. \\ \hline
\textbf{Prioridade} & Alta (Requisito legal e federativo obrigatório) \\ \hline
\textbf{Origem} & Análise Documental (Atestado Médico Desportivo) \\ \hline
\textbf{Analista} & Tiago Martins \\ \hline
\textbf{Pré-condições} & 1. Utilizador autenticado. \newline 2. Perfil base do atleta criado na BD. \\ \hline
\textbf{Pós-condições} & 1. Documento armazenado de forma segura e encriptada (RGPD). \newline 2. Estado do atleta: ``Pendente de Validação'' --- após aprovação --- ``Apto''. \newline 3. Alertas automáticos para datas de caducidade programados. \\ \hline
\textbf{Critérios de Aceitação} & 1. Formulário aceita PDF, JPG, PNG com limite de 5MB. \newline 2. Notificações a 30 e 15 dias antes de caducidade. \newline 3. Se data ultrapassada sem nova submissão válida, estado muda para ``Inapto'', bloqueando convocatórias. \\
\end{xltabular}
}
\rowcolors{2}{white}{white}

% ---------------------------------------------------------
% REQUISITO RF-12
% ---------------------------------------------------------

O requisito RF-12, denominado Gestão de Lesões e Recuperação, integra-se na área da Saúde e Bem-Estar e visa formalizar o processo de comunicação e acompanhamento de indisponibilidades físicas dos atletas. A ausência de um canal centralizado para este tipo de reporte obriga frequentemente os clubes a gerir esta informação de forma informal, com risco de envolver atletas lesionados em competições, com consequências para a saúde do desportista e para a responsabilidade do clube.

Com este módulo, o atleta ou o seu encarregado de educação pode reportar uma queixa ou lesão através de um formulário estruturado, indicando a zona corporal afetada e o nível de dor. O treinador e o departamento médico recebem uma notificação imediata e o atleta fica visualmente destacado nas listas de plantel. O sistema impede ativamente a sua inclusão em convocatórias enquanto o estado ``Lesionado'' se mantiver ativo.

{\small
\rowcolors{2}{rowblue}{white}
\begin{xltabular}{\textwidth}{|l|X|}
\hline
\rowcolor{headerblue}\textcolor{white}{\textbf{Campo}} & \textcolor{white}{\textbf{Descrição}} \\ \hline
\endfirsthead
\multicolumn{2}{c}{{\tablename\ \thetable{} -- Continuação}} \\
\hline
\rowcolor{headerblue}\textcolor{white}{\textbf{Campo}} & \textcolor{white}{\textbf{Descrição}} \\ \hline
\endhead
\hline \multicolumn{2}{|r|}{\textit{Continua na próxima página...}} \\ \hline \endfoot
\hline
\caption{Especificação do Requisito Funcional: RF-12}
\label{tab:rf-12} \\
\endlastfoot
\textbf{ID} & RF-12 \\ \hline
\textbf{Nome do Requisito} & Gestão de Lesões e Recuperação \\ \hline
\textbf{Área do Sistema} & Saúde e Bem-Estar \\ \hline
\textbf{Ator Principal} & Atletas/EE e Gerência \\ \hline
\textbf{Descrição} & O sistema deve disponibilizar um canal de registo e comunicação de lesões e indisponibilidades físicas dos atletas, permitindo ao atleta ou encarregado de educação reportar uma queixa, ao departamento médico gerir o processo de recuperação e ao treinador ter visibilidade imediata sobre a disponibilidade do plantel. \\ \hline
\textbf{Prioridade} & Alta \\ \hline
\textbf{Origem} & Entrevista --- Treinador do CAP (Reunião 2) \\ \hline
\textbf{Analista} & André Santos \\ \hline
\textbf{Pré-condições} & 1. Utilizador autenticado. \newline 2. Atleta com perfil ativo e integrado no plantel. \\ \hline
\textbf{Pós-condições} & 1. Estado do atleta atualizado para ``Lesionado'' ou ``Condicionado'' na BD. \newline 2. Notificação imediata ao treinador e fisioterapeuta. \newline 3. Atleta destacado visualmente nas listas de plantel. \\ \hline
\textbf{Critérios de Aceitação} & 1. Formulário de queixa com zona do corpo, nível de dor e descrição. \newline 2. Vista ``Zona de Lesionados'' exclusiva para Treinadores e Médicos. \newline 3. Bloqueio/alerta se treinador tenta incluir atleta ``Lesionado'' em convocatória. \\
\end{xltabular}
}
\rowcolors{2}{white}{white}

% ---------------------------------------------------------
% REQUISITO RF-13
% ---------------------------------------------------------

O requisito RF-13, designado Acompanhamento Físico e Evolução de Métricas, enquadra-se na área da Saúde e Bem-Estar e tem como finalidade suportar o acompanhamento longitudinal do desenvolvimento físico dos atletas ao longo das épocas desportivas. O registo sistemático de indicadores biométricos fornece à equipa técnica e médica informação objetiva para fundamentar decisões sobre cargas de treino, recuperação e nutrição.

O sistema agrega os dados cronologicamente e apresenta-os em gráficos de evolução filtráveis por período, permitindo identificar tendências e responder proativamente a variações preocupantes. O acesso a esta informação é estritamente confidencial, restringindo-se ao próprio atleta, ao seu encarregado de educação e aos profissionais de saúde e técnicos com ligação à equipa, em conformidade com os princípios do RGPD.

{\small
\rowcolors{2}{rowblue}{white}
\begin{xltabular}{\textwidth}{|l|X|}
\hline
\rowcolor{headerblue}\textcolor{white}{\textbf{Campo}} & \textcolor{white}{\textbf{Descrição}} \\ \hline
\endfirsthead
\multicolumn{2}{c}{{\tablename\ \thetable{} -- Continuação}} \\
\hline
\rowcolor{headerblue}\textcolor{white}{\textbf{Campo}} & \textcolor{white}{\textbf{Descrição}} \\ \hline
\endhead
\hline \multicolumn{2}{|r|}{\textit{Continua na próxima página...}} \\ \hline \endfoot
\hline
\caption{Especificação do Requisito Funcional: RF-13}
\label{tab:rf-13} \\
\endlastfoot
\textbf{ID} & RF-13 \\ \hline
\textbf{Nome do Requisito} & Acompanhamento Físico e Evolução de Métricas \\ \hline
\textbf{Área do Sistema} & Saúde e Bem-Estar \\ \hline
\textbf{Ator Principal} & Atletas/EE e Treinadores/Dept.\ Médico \\ \hline
\textbf{Descrição} & O sistema deve permitir o registo periódico e o acompanhamento longitudinal dos dados antropométricos e biométricos dos atletas (peso, altura, IMC, frequência cardíaca de repouso), gerando representações gráficas da evolução ao longo das épocas desportivas para apoio às decisões da equipa técnica e médica. \\ \hline
\textbf{Prioridade} & Média \\ \hline
\textbf{Origem} & Entrevista --- Atleta do CAP (Reunião 3) \\ \hline
\textbf{Analista} & Pedro Ferreira \\ \hline
\textbf{Pré-condições} & 1. Utilizador (treinador, médico ou atleta) autenticado. \newline 2. Perfil do atleta ativo. \\ \hline
\textbf{Pós-condições} & 1. Novos dados anexados ao processo do atleta com data de registo. \newline 2. Métricas derivadas (IMC) recalculadas e gráficos atualizados. \newline 3. Histórico disponível para equipa técnica. \\ \hline
\textbf{Critérios de Aceitação} & 1. Gráficos de linhas filtráveis por período. \newline 2. Acesso estritamente confidencial. \newline 3. Validação de dados (impede valores negativos/irrealistas). \\
\end{xltabular}
}
\rowcolors{2}{white}{white}

% ---------------------------------------------------------
% REQUISITO RF-14
% ---------------------------------------------------------

O requisito RF-14, denominado Gestão Documental e Federativa, integra-se na área da Saúde e Bem-Estar e centraliza o processo de recolha, armazenamento e validação dos documentos de identificação necessários para os processos de filiação federativa e para a ativação das apólices de seguro desportivo. A dispersão desta documentação em diferentes canais representa um risco operacional e de conformidade regulatória para o clube.

O módulo permite que atletas e encarregados de educação façam o \textit{upload} seguro dos documentos (Cartão de Cidadão, Passaporte, fotografias) diretamente na plataforma, ficando em estado de ``A aguardar validação'' até que a secretaria proceda à sua verificação. Toda a informação é armazenada de forma encriptada, em conformidade com o RGPD, e o acesso é restrito ao remetente e à gerência, sendo os treinadores expressamente impedidos de aceder a estes dados.

{\small
\rowcolors{2}{rowblue}{white}
\begin{xltabular}{\textwidth}{|l|X|}
\hline
\rowcolor{headerblue}\textcolor{white}{\textbf{Campo}} & \textcolor{white}{\textbf{Descrição}} \\ \hline
\endfirsthead
\multicolumn{2}{c}{{\tablename\ \thetable{} -- Continuação}} \\
\hline
\rowcolor{headerblue}\textcolor{white}{\textbf{Campo}} & \textcolor{white}{\textbf{Descrição}} \\ \hline
\endhead
\hline \multicolumn{2}{|r|}{\textit{Continua na próxima página...}} \\ \hline \endfoot
\hline
\caption{Especificação do Requisito Funcional: RF-14}
\label{tab:rf-14} \\
\endlastfoot
\textbf{ID} & RF-14 \\ \hline
\textbf{Nome do Requisito} & Gestão Documental e Federativa \\ \hline
\textbf{Área do Sistema} & Saúde e Bem-Estar \\ \hline
\textbf{Ator Principal} & Atletas/EE (Upload) e Gerência (Processamento) \\ \hline
\textbf{Descrição} & O sistema deve permitir o carregamento (\textit{upload}) seguro de cópias digitais de documentos de identificação oficiais (Cartão de Cidadão, Passaporte, fotos tipo passe), centralizando a informação necessária para os processos de filiação nas federações desportivas e ativação de apólices de seguro. \\ \hline
\textbf{Prioridade} & Alta (Obrigatório para competições oficiais) \\ \hline
\textbf{Origem} & Entrevista --- Gerência do CAP (Reunião 1) \\ \hline
\textbf{Analista} & Tiago Martins \\ \hline
\textbf{Pré-condições} & 1. Utilizador autenticado. \newline 2. Perfil base do atleta criado e associado a modalidade que exige inscrição federativa. \\ \hline
\textbf{Pós-condições} & 1. Ficheiros armazenados de forma encriptada (RGPD). \newline 2. Secretaria recebe alerta de novos documentos pendentes. \newline 3. Estado do atleta: ``A aguardar validação da Secretaria''. \\ \hline
\textbf{Critérios de Aceitação} & 1. Limitar ficheiros a PDF, JPG, PNG com máx. 5MB. \newline 2. Acesso restrito à Gerência e remetente (bloqueado a Treinadores e Staff). \newline 3. Secretaria pode validar ou rejeitar, gerando notificação automática. \\
\end{xltabular}
}
\rowcolors{2}{white}{white}

% ---------------------------------------------------------
% REQUISITO RF-15
% ---------------------------------------------------------

O requisito RF-15, denominado Painéis de Desempenho (\textit{Dashboards}), enquadra-se na área de Análise e Desempenho e visa disponibilizar representações visuais e dinâmicas do historial de dados desportivos acumulados na plataforma. Cada perfil de utilizador acede a uma visão adaptada ao seu nível de permissão: o atleta consulta a sua própria evolução, o treinador compara métricas do plantel e a gerência obtém uma panorâmica global do desempenho desportivo do clube.

Os painéis são gerados em modo de leitura exclusivo, com dados consolidados a partir dos módulos de assiduidade, ocorrências, lesões e métricas físicas. Filtros temporais permitem circunscrever a análise a períodos específicos (época atual, últimos 30 dias), e a gerência pode exportar os relatórios para PDF ou Excel para partilha com \textit{stakeholders} externos.

{\small
\rowcolors{2}{rowblue}{white}
\begin{xltabular}{\textwidth}{|l|X|}
\hline
\rowcolor{headerblue}\textcolor{white}{\textbf{Campo}} & \textcolor{white}{\textbf{Descrição}} \\ \hline
\endfirsthead
\multicolumn{2}{c}{{\tablename\ \thetable{} -- Continuação}} \\
\hline
\rowcolor{headerblue}\textcolor{white}{\textbf{Campo}} & \textcolor{white}{\textbf{Descrição}} \\ \hline
\endhead
\hline \multicolumn{2}{|r|}{\textit{Continua na próxima página...}} \\ \hline \endfoot
\hline
\caption{Especificação do Requisito Funcional: RF-15}
\label{tab:rf-15} \\
\endlastfoot
\textbf{ID} & RF-15 \\ \hline
\textbf{Nome do Requisito} & Painéis de Desempenho (Dashboards) \\ \hline
\textbf{Área do Sistema} & Análise e Desempenho \\ \hline
\textbf{Ator Principal} & Atletas/EE, Treinadores e Gerência \\ \hline
\textbf{Descrição} & O sistema deve gerar painéis de visualização gráfica (\textit{dashboards}) dinâmicos baseados no histórico de dados. Os atletas consultam o seu próprio desempenho, os treinadores acedem a comparações métricas do seu plantel (ex: gráficos radar) e a gerência visualiza o panorama global desportivo do clube. \\ \hline
\textbf{Prioridade} & Média \\ \hline
\textbf{Origem} & Entrevista --- Gerência do CAP (Reunião 1) \\ \hline
\textbf{Analista} & Tiago Martins \\ \hline
\textbf{Pré-condições} & 1. Utilizador autenticado. \newline 2. Sistema possui dados acumulados suficientes para gerar representações estatísticas. \\ \hline
\textbf{Pós-condições} & 1. Sistema compila dados e renderiza gráficos e tabelas ajustados ao nível de acesso. \newline 2. Dados apresentados em modo \textit{read-only}. \\ \hline
\textbf{Critérios de Aceitação} & 1. Isolamento de informação (atleta vê apenas as suas estatísticas). \newline 2. Filtros temporais (últimos 30 dias, época atual, etc.). \newline 3. Gerência pode exportar para PDF ou Excel. \\
\end{xltabular}
}
\rowcolors{2}{white}{white}

% ---------------------------------------------------------
% REQUISITO RF-16
% ---------------------------------------------------------

O requisito RF-16, designado Avaliação Qualitativa e Objetivos, integra-se na área de Análise e Desempenho e permite à equipa técnica complementar os dados quantitativos com avaliações de cariz qualitativo. O registo de notas técnicas e a definição de objetivos individuais de progressão são instrumentos fundamentais na formação desportiva, permitindo estabelecer metas claras e acompanhar a sua concretização ao longo da época.

O treinador pode redigir relatórios técnicos após treinos ou jogos e definir objetivos para cada atleta, com prazos e estados de acompanhamento. O atleta e o seu encarregado de educação têm acesso em leitura às avaliações e objetivos que lhes dizem respeito, sem visibilidade sobre os de outros elementos do plantel, garantindo total privacidade e fomento de uma cultura de \textit{feedback} individual.

{\small
\rowcolors{2}{rowblue}{white}
\begin{xltabular}{\textwidth}{|l|X|}
\hline
\rowcolor{headerblue}\textcolor{white}{\textbf{Campo}} & \textcolor{white}{\textbf{Descrição}} \\ \hline
\endfirsthead
\multicolumn{2}{c}{{\tablename\ \thetable{} -- Continuação}} \\
\hline
\rowcolor{headerblue}\textcolor{white}{\textbf{Campo}} & \textcolor{white}{\textbf{Descrição}} \\ \hline
\endhead
\hline \multicolumn{2}{|r|}{\textit{Continua na próxima página...}} \\ \hline \endfoot
\hline
\caption{Especificação do Requisito Funcional: RF-16}
\label{tab:rf-16} \\
\endlastfoot
\textbf{ID} & RF-16 \\ \hline
\textbf{Nome do Requisito} & Avaliação Qualitativa e Objetivos \\ \hline
\textbf{Área do Sistema} & Análise e Desempenho \\ \hline
\textbf{Ator Principal} & Treinadores (Registo) e Atletas/EE (Consulta) \\ \hline
\textbf{Descrição} & O sistema deve permitir à equipa técnica o registo de avaliações qualitativas (notas, relatórios técnicos) após treinos ou jogos, bem como a definição e acompanhamento de objetivos individuais de progressão para cada atleta da sua equipa. \\ \hline
\textbf{Prioridade} & Média \\ \hline
\textbf{Origem} & Entrevista --- Treinador do CAP (Reunião 2) \\ \hline
\textbf{Analista} & André Santos \\ \hline
\textbf{Pré-condições} & 1. Autenticado com perfil de Treinador. \newline 2. Atleta a avaliar inscrito e integrado no plantel. \\ \hline
\textbf{Pós-condições} & 1. Relatório/objetivo persistido e anexado ao processo desportivo. \newline 2. Atleta/EE tem acesso de leitura na área pessoal. \newline 3. Notificação informando de nova avaliação/objetivo. \\ \hline
\textbf{Critérios de Aceitação} & 1. Formulário com pontuações rápidas (escalas/estrelas) e texto livre. \newline 2. Objetivos com \textit{deadlines} e estados (``Em curso'', ``Atingido'', ``Não Atingido''). \newline 3. Total privacidade (atleta não vê avaliações de colegas). \\
\end{xltabular}
}
\rowcolors{2}{white}{white}

% ---------------------------------------------------------
% REQUISITO RF-17
% ---------------------------------------------------------

O requisito RF-17, denominado Gestão de Quotas e Pagamentos, integra-se na área da Tesouraria e Contabilidade e representa um pilar crítico para a sustentabilidade financeira do clube. A automatização do ciclo de cobranças -- desde a geração de referências de pagamento até à emissão de recibos e ao controlo de dívidas -- reduz significativamente a carga administrativa da tesouraria e melhora a taxa de recebimento de quotas.

O sistema integra-se com \textit{gateways} de pagamento externos (como SIBS/Multibanco e MBWay) para disponibilizar referências de pagamento automáticas. As transações são averbadas automaticamente na conta corrente do associado, e a fatura/recibo é gerada em PDF para \textit{download} imediato. Em caso de dívida superior a dois meses, o sistema aplica bloqueios desportivos de forma automática, garantindo que a gestão de morosidade é sistemática e consistente.

{\small
\rowcolors{2}{rowblue}{white}
\begin{xltabular}{\textwidth}{|l|X|}
\hline
\rowcolor{headerblue}\textcolor{white}{\textbf{Campo}} & \textcolor{white}{\textbf{Descrição}} \\ \hline
\endfirsthead
\multicolumn{2}{c}{{\tablename\ \thetable{} -- Continuação}} \\
\hline
\rowcolor{headerblue}\textcolor{white}{\textbf{Campo}} & \textcolor{white}{\textbf{Descrição}} \\ \hline
\endhead
\hline \multicolumn{2}{|r|}{\textit{Continua na próxima página...}} \\ \hline \endfoot
\hline
\caption{Especificação do Requisito Funcional: RF-17}
\label{tab:rf-17} \\
\endlastfoot
\textbf{ID} & RF-17 \\ \hline
\textbf{Nome do Requisito} & Gestão de Quotas e Pagamentos \\ \hline
\textbf{Área do Sistema} & Tesouraria e Contabilidade \\ \hline
\textbf{Ator Principal} & Atletas/Encarregados de Educação (Pagamento) e Tesouraria/Gerência (Gestão) \\ \hline
\textbf{Descrição} & O sistema deve incluir um motor financeiro integrado para a gestão de contas correntes, permitindo a geração de referências de pagamento automáticas (Multibanco, MBWay), a emissão de recibos e a monitorização de valores em atraso. \\ \hline
\textbf{Prioridade} & Alta (Crítico para a sustentabilidade do clube) \\ \hline
\textbf{Origem} & Análise Documental (Fatura/Recibo de Quota) \\ \hline
\textbf{Analista} & Pedro Ferreira \\ \hline
\textbf{Pré-condições} & 1. Utilizador autenticado. \newline 2. Integração com \textit{gateways} de pagamento externos (SIBS, Ifthenpay) configurada e ativa. \newline 3. Plano de pagamentos previamente definido pela gerência. \\ \hline
\textbf{Pós-condições} & 1. Transação averbada automaticamente na conta corrente do utilizador. \newline 2. Fatura/recibo em PDF disponibilizado para \textit{download}. \newline 3. Suspensões desportivas levantadas se dívida regularizada. \\ \hline
\textbf{Critérios de Aceitação} & 1. Perfil parental consolida quotas de vários educandos. \newline 2. Bloqueio automático (mais de 2 meses de quotas em atraso). \newline 3. Vista ``Devedores'' em tempo real com motor de notificações. \\
\end{xltabular}
}
\rowcolors{2}{white}{white}

% ---------------------------------------------------------
% REQUISITO RF-18
% ---------------------------------------------------------

O requisito RF-18, designado Contabilidade Global do Clube, enquadra-se na área da Tesouraria e Contabilidade e visa disponibilizar à gerência e à tesouraria uma visão centralizada e em tempo real de toda a atividade financeira da organização. Para além das quotas dos associados, o clube gere um conjunto vasto de despesas e receitas -- contratos de \textit{staff} técnico, patrocínios, despesas com fornecedores de serviços essenciais -- que exigem uma plataforma de controlo integrada.

O módulo agrega todos os movimentos financeiros registados e apresenta-os em painéis com gráficos comparativos de rendimentos versus gastos, desagregados por categoria. O registo de faturas de fornecedores inclui alertas de data limite de pagamento, e a exportação dos dados para PDF ou Excel facilita a prestação de contas a órgãos sociais ou entidades externas. O acesso a este módulo é estritamente restrito aos perfis de topo da organização.

{\small
\rowcolors{2}{rowblue}{white}
\begin{xltabular}{\textwidth}{|l|X|}
\hline
\rowcolor{headerblue}\textcolor{white}{\textbf{Campo}} & \textcolor{white}{\textbf{Descrição}} \\ \hline
\endfirsthead
\multicolumn{2}{c}{{\tablename\ \thetable{} -- Continuação}} \\
\hline
\rowcolor{headerblue}\textcolor{white}{\textbf{Campo}} & \textcolor{white}{\textbf{Descrição}} \\ \hline
\endhead
\hline \multicolumn{2}{|r|}{\textit{Continua na próxima página...}} \\ \hline \endfoot
\hline
\caption{Especificação do Requisito Funcional: RF-18}
\label{tab:rf-18} \\
\endlastfoot
\textbf{ID} & RF-18 \\ \hline
\textbf{Nome do Requisito} & Contabilidade Global do Clube \\ \hline
\textbf{Área do Sistema} & Tesouraria e Contabilidade \\ \hline
\textbf{Ator Principal} & Gerência/Tesouraria \\ \hline
\textbf{Descrição} & O sistema deve disponibilizar painéis de controlo (\textit{dashboards}) financeiros avançados para a gerência, permitindo a gestão e visualização centralizada de todas as despesas e receitas, incluindo contratos (vencimentos da equipa técnica e \textit{staff}), patrocínios e pagamentos a fornecedores (água, eletricidade, transportes). \\ \hline
\textbf{Prioridade} & Alta \\ \hline
\textbf{Origem} & Entrevista --- Gerência do CAP (Reunião 1) \\ \hline
\textbf{Analista} & Tiago Martins \\ \hline
\textbf{Pré-condições} & 1. Autenticado com perfil restrito de topo (Gerência, Tesouraria ou Administração). \newline 2. Movimentos financeiros registados na BD. \\ \hline
\textbf{Pós-condições} & 1. Sistema agrega dados e gera relatórios de balanço em tempo real. \newline 2. Dados exportáveis (PDF ou Excel). \\ \hline
\textbf{Critérios de Aceitação} & 1. Gráficos circulares/barras comparando rendimentos vs. gastos. \newline 2. Registo de faturas de fornecedores com alertas de data limite. \newline 3. Acesso estritamente bloqueado a Atletas, EE e Treinadores. \\
\end{xltabular}
}
\rowcolors{2}{white}{white}

% ---------------------------------------------------------
% REQUISITO RF-19
% ---------------------------------------------------------

O requisito RF-19, denominado Conformidade Fiscal (Exportação SAF-T), insere-se na área da Tesouraria e Contabilidade e trata do cumprimento das obrigações fiscais do clube perante a Autoridade Tributária e Aduaneira (AT). A comunicação das receitas através de ficheiros normalizados (Standard Audit File for Tax -- SAF-T) é uma obrigação legal que, quando não automatizada, representa um processo manual moroso e sujeito a erros humanos.

O sistema compila automaticamente as transações validadas num intervalo temporal selecionável e gera o ficheiro XML no formato SAF-T conforme as especificações das Finanças. Antes da geração final, é apresentado um relatório de validação que identifica eventuais anomalias, garantindo que o ficheiro submetido está em conformidade. O acesso a esta funcionalidade é restrito aos perfis financeiros de topo.

{\small
\rowcolors{2}{rowblue}{white}
\begin{xltabular}{\textwidth}{|l|X|}
\hline
\rowcolor{headerblue}\textcolor{white}{\textbf{Campo}} & \textcolor{white}{\textbf{Descrição}} \\ \hline
\endfirsthead
\multicolumn{2}{c}{{\tablename\ \thetable{} -- Continuação}} \\
\hline
\rowcolor{headerblue}\textcolor{white}{\textbf{Campo}} & \textcolor{white}{\textbf{Descrição}} \\ \hline
\endhead
\hline \multicolumn{2}{|r|}{\textit{Continua na próxima página...}} \\ \hline \endfoot
\hline
\caption{Especificação do Requisito Funcional: RF-19}
\label{tab:rf-19} \\
\endlastfoot
\textbf{ID} & RF-19 \\ \hline
\textbf{Nome do Requisito} & Conformidade Fiscal (Exportação SAF-T) \\ \hline
\textbf{Área do Sistema} & Tesouraria e Contabilidade \\ \hline
\textbf{Ator Principal} & Gerência/Tesouraria/Contabilista \\ \hline
\textbf{Descrição} & O sistema deve garantir o cumprimento das obrigações legais através da geração e exportação de ficheiros de faturação normalizados (como o formato SAF-T), permitindo a comunicação das receitas do clube (quotas, \textit{merchandising}, alugueres) à Autoridade Tributária. \\ \hline
\textbf{Prioridade} & Alta (Obrigatório por lei) \\ \hline
\textbf{Origem} & Entrevista --- Gerência do CAP (Reunião 1) \\ \hline
\textbf{Analista} & Tiago Martins \\ \hline
\textbf{Pré-condições} & 1. Autenticado com perfil de topo financeiro. \newline 2. Dados fiscais do clube configurados (NIF, morada, série de faturação). \newline 3. Recibos/faturas emitidos e validados no período a exportar. \\ \hline
\textbf{Pós-condições} & 1. Sistema compila transações e gera ficheiro XML conforme Finanças. \newline 2. Ficheiro disponível para \textit{download} seguro. \newline 3. Registo de auditoria interno da extração. \\ \hline
\textbf{Critérios de Aceitação} & 1. Seleção precisa do intervalo temporal. \newline 2. Validação prévia com relatório de erros (anomalias). \newline 3. Acesso estritamente restrito a perfis financeiros autorizados. \\
\end{xltabular}
}
\rowcolors{2}{white}{white}

% ---------------------------------------------------------
% REQUISITO RF-20
% ---------------------------------------------------------

O requisito RF-20, denominado Gestão de Espaços Físicos, integra-se na área de Logística, Infraestruturas e Equipamentos e resolve o problema recorrente da gestão ineficaz da utilização dos espaços do clube. A ausência de um sistema centralizado de reservas conduz frequentemente a conflitos de agendamento, subutilização de instalações e dificuldades no controlo de alugueres a entidades externas.

O módulo disponibiliza um calendário unificado onde treinadores e gerência podem reservar campos, pavilhões e outros espaços para as suas atividades. O sistema valida proativamente a disponibilidade e impede marcações em duplicado. Para alugueres a entidades externas, o sistema aciona automaticamente a notificação da tesouraria para a emissão da respetiva fatura, integrando a gestão de espaços com o módulo financeiro.

{\small
\rowcolors{2}{rowblue}{white}
\begin{xltabular}{\textwidth}{|l|X|}
\hline
\rowcolor{headerblue}\textcolor{white}{\textbf{Campo}} & \textcolor{white}{\textbf{Descrição}} \\ \hline
\endfirsthead
\multicolumn{2}{c}{{\tablename\ \thetable{} -- Continuação}} \\
\hline
\rowcolor{headerblue}\textcolor{white}{\textbf{Campo}} & \textcolor{white}{\textbf{Descrição}} \\ \hline
\endhead
\hline \multicolumn{2}{|r|}{\textit{Continua na próxima página...}} \\ \hline \endfoot
\hline
\caption{Especificação do Requisito Funcional: RF-20}
\label{tab:rf-20} \\
\endlastfoot
\textbf{ID} & RF-20 \\ \hline
\textbf{Nome do Requisito} & Gestão de Espaços Físicos \\ \hline
\textbf{Área do Sistema} & Logística, Infraestruturas e Equipamentos \\ \hline
\textbf{Ator Principal} & Gerência e Treinadores \\ \hline
\textbf{Descrição} & O sistema deve disponibilizar um calendário unificado para a gestão e reserva dos espaços físicos do clube (campos, pavilhões, balneários, ginásios), garantindo a alocação correta para os treinos/jogos e permitindo o registo e controlo de alugueres a entidades externas. \\ \hline
\textbf{Prioridade} & Alta \\ \hline
\textbf{Origem} & Entrevista --- Gerência do CAP (Reunião 1) \\ \hline
\textbf{Analista} & Tiago Martins \\ \hline
\textbf{Pré-condições} & 1. Autenticado com permissões de gestão de horários (Treinador ou Gerência). \newline 2. Espaços físicos e características parametrizados no sistema. \\ \hline
\textbf{Pós-condições} & 1. Reserva registada e refletida imediatamente no calendário global. \newline 2. Espaço vinculado à sessão da equipa. \newline 3. Notificação à tesouraria para faturação de aluguer externo. \\ \hline
\textbf{Critérios de Aceitação} & 1. Bloqueio proativo de dupla marcação. \newline 2. Calendário com vista diária/semanal/mensal com códigos de cores. \newline 3. Gerência define períodos de inatividade (obras/manutenção). \\
\end{xltabular}
}
\rowcolors{2}{white}{white}

% ---------------------------------------------------------
% REQUISITO RF-21
% ---------------------------------------------------------

O requisito RF-21, designado Manutenção de Infraestruturas, enquadra-se na área de Logística, Infraestruturas e Equipamentos e visa organizar e registar as intervenções de manutenção nas instalações do clube. A manutenção preventiva e curativa de campos relvados, sistemas de iluminação e balneários é essencial para garantir condições de treino adequadas e para preservar o valor patrimonial das infraestruturas.

O sistema permite agendar intervenções e registar o seu estado (``Agendada'', ``Em curso'', ``Concluída''), bloqueando automaticamente novas reservas no espaço afetado durante o período de indisponibilidade. Os treinadores e equipas com atividades planeadas no espaço são notificados proativamente, e os custos associados são registados e integrados na contabilidade do clube.

{\small
\rowcolors{2}{rowblue}{white}
\begin{xltabular}{\textwidth}{|l|X|}
\hline
\rowcolor{headerblue}\textcolor{white}{\textbf{Campo}} & \textcolor{white}{\textbf{Descrição}} \\ \hline
\endfirsthead
\multicolumn{2}{c}{{\tablename\ \thetable{} -- Continuação}} \\
\hline
\rowcolor{headerblue}\textcolor{white}{\textbf{Campo}} & \textcolor{white}{\textbf{Descrição}} \\ \hline
\endhead
\hline \multicolumn{2}{|r|}{\textit{Continua na próxima página...}} \\ \hline \endfoot
\hline
\caption{Especificação do Requisito Funcional: RF-21}
\label{tab:rf-21} \\
\endlastfoot
\textbf{ID} & RF-21 \\ \hline
\textbf{Nome do Requisito} & Manutenção de Infraestruturas \\ \hline
\textbf{Área do Sistema} & Logística, Infraestruturas e Equipamentos \\ \hline
\textbf{Ator Principal} & Gerência/Staff (Manutenção) \\ \hline
\textbf{Descrição} & O sistema deve permitir o agendamento e o registo de intervenções de manutenção preventivas e curativas nas infraestruturas do clube (ex: tratamento de relvados, afinação de iluminação, reparação de balneários), articulando essas ações com a disponibilidade do calendário. \\ \hline
\textbf{Prioridade} & Média \\ \hline
\textbf{Origem} & Entrevista --- Gerência do CAP (Reunião 1) \\ \hline
\textbf{Analista} & Tiago Martins \\ \hline
\textbf{Pré-condições} & 1. Autenticado com perfil Gerência ou \textit{Staff} de manutenção. \newline 2. Espaços físicos parametrizados no sistema. \\ \hline
\textbf{Pós-condições} & 1. Registo guardado no histórico do espaço físico. \newline 2. Bloqueio automático de novas reservas durante a manutenção. \newline 3. Notificação a treinadores/equipas afetados. \\ \hline
\textbf{Critérios de Aceitação} & 1. Estados de intervenção (``Agendada'', ``Em curso'', ``Concluída''). \newline 2. Integração com contabilidade para registo de custos (faturas/orçamentos). \newline 3. Bloqueio visualmente distinto no calendário. \\
\end{xltabular}
}
\rowcolors{2}{white}{white}

% ---------------------------------------------------------
% REQUISITO RF-22
% ---------------------------------------------------------

O requisito RF-22, denominado Gestão de Stock e Loja, integra-se na área de Logística, Infraestruturas e Equipamentos e tem como objetivo centralizar a gestão do inventário de material desportivo e do fardamento do clube. A gestão manual de stock origina frequentemente ruturas de material durante a época e dificulta o controlo de encomendas individuais de fardamento por parte dos atletas.

O módulo disponibiliza uma loja interna na plataforma onde atletas e encarregados de educação podem encomendar e pagar fardamento, com o movimento financeiro a ser automaticamente gerado na conta corrente do associado. Em paralelo, a gerência e os treinadores dispõem de ferramentas de controlo de stock, com alertas de nível mínimo e distinção clara entre ``Material de Venda'' e ``Material de Treino'' para uso coletivo.

{\small
\rowcolors{2}{rowblue}{white}
\begin{xltabular}{\textwidth}{|l|X|}
\hline
\rowcolor{headerblue}\textcolor{white}{\textbf{Campo}} & \textcolor{white}{\textbf{Descrição}} \\ \hline
\endfirsthead
\multicolumn{2}{c}{{\tablename\ \thetable{} -- Continuação}} \\
\hline
\rowcolor{headerblue}\textcolor{white}{\textbf{Campo}} & \textcolor{white}{\textbf{Descrição}} \\ \hline
\endhead
\hline \multicolumn{2}{|r|}{\textit{Continua na próxima página...}} \\ \hline \endfoot
\hline
\caption{Especificação do Requisito Funcional: RF-22}
\label{tab:rf-22} \\
\endlastfoot
\textbf{ID} & RF-22 \\ \hline
\textbf{Nome do Requisito} & Gestão de Stock e Loja \\ \hline
\textbf{Área do Sistema} & Logística, Infraestruturas e Equipamentos \\ \hline
\textbf{Ator Principal} & Gerência/Treinadores (Controlo) e Atletas/Encarregados de Educação (Encomendas) \\ \hline
\textbf{Descrição} & O sistema deve disponibilizar um módulo de loja interna e gestão de inventário, permitindo aos atletas encomendar e pagar fardamento através da plataforma. Simultaneamente, permite à gerência e aos treinadores controlar o \textit{stock} existente e gerir a requisição de material desportivo (bolas, coletes, cones) para o dia a dia. \\ \hline
\textbf{Prioridade} & Média \\ \hline
\textbf{Origem} & Questionário Online (47 respostas) \\ \hline
\textbf{Analista} & Pedro Ferreira \\ \hline
\textbf{Pré-condições} & 1. Utilizador autenticado com o perfil adequado. \newline 2. Catálogo de produtos e fardamento parametrizado no sistema com preços e níveis de \textit{stock} definidos. \\ \hline
\textbf{Pós-condições} & 1. Encomenda do atleta gera movimento na conta corrente e referência de pagamento/recibo. \newline 2. \textit{Stock} atualizado (deduzido). \newline 3. Gerência/rouparia notificada. \\ \hline
\textbf{Critérios de Aceitação} & 1. Interface da loja intuitiva com tamanhos e preço final. \newline 2. Alerta de \textit{stock} mínimo. \newline 3. Distinção entre ``Material de Venda'' e ``Material de Treino''. \\
\end{xltabular}
}
\rowcolors{2}{white}{white}

% ---------------------------------------------------------
% REQUISITO RF-23
% ---------------------------------------------------------

O requisito RF-23, designado Gestão de Eventos Extra-desportivos, enquadra-se na área de Logística, Infraestruturas e Equipamentos e visa estender as capacidades de organização da plataforma a eventos de natureza social e comunitária promovidos pelo clube. Galas de fim de época, jantares temáticos e torneios solidários são momentos de coesão associativa que exigem uma logística de inscrições e cobranças que, quando gerida manualmente, consome recursos significativos da secretaria.

A gerência pode criar eventos, definir lotação máxima, prazo de inscrição, público-alvo e eventual custo de ingresso. Os membros inscrevem-se diretamente na plataforma, sendo gerada automaticamente uma referência de pagamento quando aplicável. O sistema bloqueia novas inscrições quando a lotação ou o prazo são atingidos, e disponibiliza uma lista de inscritos exportável para apoiar a organização do evento.

{\small
\rowcolors{2}{rowblue}{white}
\begin{xltabular}{\textwidth}{|l|X|}
\hline
\rowcolor{headerblue}\textcolor{white}{\textbf{Campo}} & \textcolor{white}{\textbf{Descrição}} \\ \hline
\endfirsthead
\multicolumn{2}{c}{{\tablename\ \thetable{} -- Continuação}} \\
\hline
\rowcolor{headerblue}\textcolor{white}{\textbf{Campo}} & \textcolor{white}{\textbf{Descrição}} \\ \hline
\endhead
\hline \multicolumn{2}{|r|}{\textit{Continua na próxima página...}} \\ \hline \endfoot
\hline
\caption{Especificação do Requisito Funcional: RF-23}
\label{tab:rf-23} \\
\endlastfoot
\textbf{ID} & RF-23 \\ \hline
\textbf{Nome do Requisito} & Gestão de Eventos Extra-desportivos \\ \hline
\textbf{Área do Sistema} & Logística, Infraestruturas e Equipamentos \\ \hline
\textbf{Ator Principal} & Gerência (Organização) e Todos os Utilizadores (Inscrição) \\ \hline
\textbf{Descrição} & O sistema deve permitir a criação, divulgação e gestão de inscrições para eventos extra-desportivos promovidos pelo clube (ex: galas de fim de época, jantares de Natal, torneios solidários), centralizando a confirmação de presenças e a cobrança de eventuais ingressos. \\ \hline
\textbf{Prioridade} & Baixa (Complementar) \\ \hline
\textbf{Origem} & Questionário Online (47 respostas) \\ \hline
\textbf{Analista} & Pedro Ferreira \\ \hline
\textbf{Pré-condições} & 1. Autenticado com perfil de Gerência. \newline 2. Módulos de Notificações (RF-04) e Pagamentos (RF-17) operacionais. \\ \hline
\textbf{Pós-condições} & 1. Evento visível no calendário global para públicos-alvo selecionados. \newline 2. Inscrição gera adição à lista de presenças e referência de pagamento. \newline 3. Painel de controlo em tempo real com inscritos e receitas. \\ \hline
\textbf{Critérios de Aceitação} & 1. Formulário define lotação, data limite, público-alvo e custo. \newline 2. Bloqueio automático quando lotação/data atingida. \newline 3. Exportação da lista de inscritos (PDF ou Excel). \\
\end{xltabular}
}
\rowcolors{2}{white}{white}

\section{Requisitos Não-Funcionais}\label{sec:requisitos_nao_funcionais}

\noindent\textit{Atalho: Saltar para o \hyperref[chap:especificacao_modelacao]{Próximo Capítulo (Modelação)}.}\par\medskip

% ---------------------------------------------------------
% REQUISITO RNF-01
% ---------------------------------------------------------

O requisito não funcional RNF-01, relativo à Proteção de Dados Sensíveis através de Encriptação, enquadra-se na categoria de Segurança e Privacidade (RGPD) e estabelece as normas técnicas de proteção da informação mais sensível gerida pela plataforma. O CAP armazena dados de saúde, documentos de identificação e informação financeira de menores e adultos, pelo que a proteção rigorosa destes dados é simultaneamente uma obrigação legal e uma condição de confiança dos associados.

O sistema deve garantir que toda a comunicação entre clientes e servidores é realizada exclusivamente sobre HTTPS, que as palavras-passe são armazenadas com algoritmos de \textit{hashing} resistentes a ataques de força bruta (como bcrypt ou Argon2) e que os ficheiros de saúde e identificação são encriptados no servidor de armazenamento antes de persistidos.

{\small
\rowcolors{2}{rowblue}{white}
\begin{xltabular}{\textwidth}{|l|X|}
\hline
\rowcolor{headerblue}\textcolor{white}{\textbf{Campo}} & \textcolor{white}{\textbf{Descrição}} \\ \hline
\endfirsthead
\multicolumn{2}{c}{{\tablename\ \thetable{} -- Continuação}} \\
\hline
\rowcolor{headerblue}\textcolor{white}{\textbf{Campo}} & \textcolor{white}{\textbf{Descrição}} \\ \hline
\endhead
\hline \multicolumn{2}{|r|}{\textit{Continua na próxima página...}} \\ \hline \endfoot
\hline
\caption{Especificação do Requisito Não Funcional: RNF-01}
\label{tab:rnf-01} \\
\endlastfoot
\textbf{ID} & RNF-01 \\ \hline
\textbf{Nome do Requisito} & Proteção de Dados Sensíveis (Encriptação) \\ \hline
\textbf{Categoria} & Segurança e Privacidade (RGPD) \\ \hline
\textbf{Descrição} & O sistema deve aplicar encriptação rigorosa a toda a informação sensível, tanto em repouso como em trânsito, protegendo dados de saúde, documentos de identificação e credenciais de acesso de todos os utilizadores da plataforma. \\ \hline
\textbf{Prioridade} & Alta (Crítico/Legal) \\ \hline
\textbf{Origem} & Equipa de Desenvolvimento \\ \hline
\textbf{Analista} & Daniel Parente \\ \hline
\textbf{Critério de Verificação} & HTTPS forçado em todas as comunicações; palavras-passe armazenadas com \textit{hashing} forte (bcrypt ou Argon2); ficheiros de saúde e identificação encriptados no \textit{storage}. \\
\end{xltabular}
}
\rowcolors{2}{white}{white}

% ---------------------------------------------------------
% REQUISITO RNF-02
% ---------------------------------------------------------

O requisito RNF-02, denominado Conformidade Legal e RGPD, integra-se na categoria de Segurança e Privacidade e estabelece os princípios de tratamento de dados pessoais que a plataforma deve respeitar ao abrigo do Regulamento Geral sobre a Proteção de Dados. A natureza do sistema -- que envolve dados de menores, dados de saúde e dados financeiros -- torna o cumprimento do RGPD não apenas obrigatório por lei, mas também essencial para a confiança dos associados na organização.

A plataforma deve implementar mecanismos de recolha de consentimento explícito e informado, garantir ao titular o direito de acesso, portabilidade e supressão (direito ao esquecimento) dos seus dados, e assegurar que o tratamento de dados de menores é devidamente autorizado. O sistema deve também garantir que, em caso de incidente de segurança, os procedimentos de notificação previstos no RGPD possam ser ativados com base nos registos de auditoria disponíveis.

{\small
\rowcolors{2}{rowblue}{white}
\begin{xltabular}{\textwidth}{|l|X|}
\hline
\rowcolor{headerblue}\textcolor{white}{\textbf{Campo}} & \textcolor{white}{\textbf{Descrição}} \\ \hline
\endfirsthead
\multicolumn{2}{c}{{\tablename\ \thetable{} -- Continuação}} \\
\hline
\rowcolor{headerblue}\textcolor{white}{\textbf{Campo}} & \textcolor{white}{\textbf{Descrição}} \\ \hline
\endhead
\hline \multicolumn{2}{|r|}{\textit{Continua na próxima página...}} \\ \hline \endfoot
\hline
\caption{Especificação do Requisito Não Funcional: RNF-02}
\label{tab:rnf-02} \\
\endlastfoot
\textbf{ID} & RNF-02 \\ \hline
\textbf{Nome do Requisito} & Conformidade Legal e RGPD \\ \hline
\textbf{Categoria} & Segurança e Privacidade (RGPD) \\ \hline
\textbf{Descrição} & O sistema deve cumprir integralmente o RGPD, implementando gestão de consentimentos, o direito ao esquecimento e proteção reforçada para dados de menores, garantindo que o tratamento de dados pessoais é lícito, transparente e limitado à finalidade declarada. \\ \hline
\textbf{Prioridade} & Alta (Crítico/Legal) \\ \hline
\textbf{Origem} & Equipa de Desenvolvimento + Entrevista Gerência \\ \hline
\textbf{Analista} & Tiago Martins \\ \hline
\textbf{Critério de Verificação} & \textit{Checkboxes} de consentimento não pré-preenchidas; exportação estruturada de dados (portabilidade); anonimização irreversível disponível mediante pedido do titular. \\
\end{xltabular}
}
\rowcolors{2}{white}{white}

% ---------------------------------------------------------
% REQUISITO RNF-03
% ---------------------------------------------------------

O requisito RNF-03, denominado Controlo de Acessos (RBAC), enquadra-se na categoria de Segurança e Privacidade e define a arquitetura de autorização da plataforma com base no modelo de controlo de acesso por papéis (\textit{Role-Based Access Control}). Num sistema com múltiplos tipos de utilizadores (Atletas, Encarregados de Educação, Treinadores, \textit{Staff}, Gerência), a segregação rigorosa das permissões é fundamental para garantir que cada perfil acede exclusivamente à informação e às funcionalidades que lhe são pertinentes.

O modelo RBAC deve ser implementado em todas as camadas da aplicação: as permissões verificadas no lado do servidor impedem acessos não autorizados através de chamadas diretas à API, enquanto a interface oculta dinamicamente menus, botões e secções para as quais o utilizador não tem autorização. A nível horizontal, o isolamento entre equipas garante que um treinador de uma modalidade não consegue aceder a informação de outra equipa, mesmo partilhando o mesmo nível hierárquico.

{\small
\rowcolors{2}{rowblue}{white}
\begin{xltabular}{\textwidth}{|l|X|}
\hline
\rowcolor{headerblue}\textcolor{white}{\textbf{Campo}} & \textcolor{white}{\textbf{Descrição}} \\ \hline
\endfirsthead
\multicolumn{2}{c}{{\tablename\ \thetable{} -- Continuação}} \\
\hline
\rowcolor{headerblue}\textcolor{white}{\textbf{Campo}} & \textcolor{white}{\textbf{Descrição}} \\ \hline
\endhead
\hline \multicolumn{2}{|r|}{\textit{Continua na próxima página...}} \\ \hline \endfoot
\hline
\caption{Especificação do Requisito Não Funcional: RNF-03}
\label{tab:rnf-03} \\
\endlastfoot
\textbf{ID} & RNF-03 \\ \hline
\textbf{Nome do Requisito} & Controlo de Acessos (RBAC) \\ \hline
\textbf{Categoria} & Segurança e Privacidade (RGPD) \\ \hline
\textbf{Descrição} & O sistema deve implementar um modelo de controlo de acesso baseado em papéis (RBAC) com privilégios explícitos por cargo, garantindo que cada utilizador acede apenas às funcionalidades e dados correspondentes ao seu perfil, tanto ao nível vertical (hierarquia) como horizontal (equipa/modalidade). \\ \hline
\textbf{Prioridade} & Alta \\ \hline
\textbf{Origem} & Entrevista --- Gerência do CAP (Reunião 1) \\ \hline
\textbf{Analista} & Tiago Martins \\ \hline
\textbf{Critério de Verificação} & Resposta HTTP 403 para qualquer acesso sem permissões via API; interface oculta menus e botões sem autorização; isolamento horizontal por equipa verificável em testes de penetração. \\
\end{xltabular}
}
\rowcolors{2}{white}{white}

% ---------------------------------------------------------
% REQUISITO RNF-04
% ---------------------------------------------------------

O requisito RNF-04, denominado Design Responsivo (Mobile-First), enquadra-se na categoria de Usabilidade e Experiência de Utilização e estabelece que a interface da plataforma deve ser concebida prioritariamente para dispositivos móveis, adaptando-se fluidamente a ecrãs de diferentes dimensões. A natureza operacional de muitos utilizadores da plataforma -- treinadores a realizar chamadas no campo, atletas a confirmar convocatórias em movimento -- impõe que a experiência em dispositivos móveis seja de primeira qualidade.

A abordagem \textit{mobile-first} implica que os \textit{layouts}, tipografias e elementos de interação são desenhados primeiramente para ecrãs pequenos e progressivamente adaptados para ecrãs maiores. Todos os elementos interativos devem ter dimensões adequadas para interação por toque, e a reorganização dos componentes deve ser fluida e intuitiva independentemente do dispositivo utilizado.

{\small
\rowcolors{2}{rowblue}{white}
\begin{xltabular}{\textwidth}{|l|X|}
\hline
\rowcolor{headerblue}\textcolor{white}{\textbf{Campo}} & \textcolor{white}{\textbf{Descrição}} \\ \hline
\endfirsthead
\multicolumn{2}{c}{{\tablename\ \thetable{} -- Continuação}} \\
\hline
\rowcolor{headerblue}\textcolor{white}{\textbf{Campo}} & \textcolor{white}{\textbf{Descrição}} \\ \hline
\endhead
\hline \multicolumn{2}{|r|}{\textit{Continua na próxima página...}} \\ \hline \endfoot
\hline
\caption{Especificação do Requisito Não Funcional: RNF-04}
\label{tab:rnf-04} \\
\endlastfoot
\textbf{ID} & RNF-04 \\ \hline
\textbf{Nome do Requisito} & Design Responsivo (Mobile-First) \\ \hline
\textbf{Categoria} & Usabilidade e Experiência de Utilização \\ \hline
\textbf{Descrição} & A interface da plataforma deve ser desenvolvida segundo uma abordagem \textit{mobile-first}, garantindo que todos os ecrãs e funcionalidades se reorganizam e adaptam de forma fluida a dispositivos de diferentes dimensões, desde smartphones até monitores de secretária. \\ \hline
\textbf{Prioridade} & Média \\ \hline
\textbf{Origem} & Questionário Online (47 respostas) \\ \hline
\textbf{Analista} & Pedro Ferreira \\ \hline
\textbf{Critério de Verificação} & A interface reorganiza-se corretamente em ecrãs pequenos (320px--768px); elementos de interação com dimensões mínimas adequadas para toque (48x48dp); ausência de \textit{scroll} horizontal não intencional em dispositivos móveis. \\
\end{xltabular}
}
\rowcolors{2}{white}{white}

% ---------------------------------------------------------
% REQUISITO RNF-05
% ---------------------------------------------------------

O requisito RNF-05, denominado Interface Intuitiva e Acessibilidade, integra-se na categoria de Usabilidade e Experiência de Utilização e visa garantir que a plataforma é utilizável por um público com diferentes níveis de literacia digital, desde jovens atletas até encarregados de educação com menor familiaridade com tecnologia. A adoção da plataforma pelo clube depende diretamente da facilidade com que os utilizadores conseguem completar as tarefas sem necessidade de formação extensiva.

O sistema deve proporcionar \textit{feedback} imediato após cada ação (confirmações, alertas de erro, notificações de sucesso), garantir que qualquer funcionalidade é alcançável em no máximo três cliques a partir do ecrã principal e respeitar os critérios de acessibilidade WCAG relativamente ao contraste de cores e à legibilidade dos elementos de interface.

{\small
\rowcolors{2}{rowblue}{white}
\begin{xltabular}{\textwidth}{|l|X|}
\hline
\rowcolor{headerblue}\textcolor{white}{\textbf{Campo}} & \textcolor{white}{\textbf{Descrição}} \\ \hline
\endfirsthead
\multicolumn{2}{c}{{\tablename\ \thetable{} -- Continuação}} \\
\hline
\rowcolor{headerblue}\textcolor{white}{\textbf{Campo}} & \textcolor{white}{\textbf{Descrição}} \\ \hline
\endhead
\hline \multicolumn{2}{|r|}{\textit{Continua na próxima página...}} \\ \hline \endfoot
\hline
\caption{Especificação do Requisito Não Funcional: RNF-05}
\label{tab:rnf-05} \\
\endlastfoot
\textbf{ID} & RNF-05 \\ \hline
\textbf{Nome do Requisito} & Interface Intuitiva e Acessibilidade \\ \hline
\textbf{Categoria} & Usabilidade e Experiência de Utilização \\ \hline
\textbf{Descrição} & A plataforma deve apresentar uma curva de aprendizagem reduzida, com navegação clara e \textit{feedback} imediato após cada interação, sendo inclusiva para utilizadores com diferentes níveis de literacia digital. A interface deve cumprir os critérios de acessibilidade mínimos definidos pelas normas WCAG. \\ \hline
\textbf{Prioridade} & Média \\ \hline
\textbf{Origem} & Questionário Online (47 respostas) \\ \hline
\textbf{Analista} & Pedro Ferreira \\ \hline
\textbf{Critério de Verificação} & \textit{Toasts}/alertas de \textit{feedback} imediatos após cada ação; qualquer funcionalidade principal acessível em no máximo 3 cliques; rácio de contraste de cores conforme WCAG 2.1 nível AA. \\
\end{xltabular}
}
\rowcolors{2}{white}{white}

% ---------------------------------------------------------
% REQUISITO RNF-06
% ---------------------------------------------------------

O requisito RNF-06, relativo aos Tempos de Resposta Rápidos (\textit{Performance}), enquadra-se na categoria de Desempenho e Escalabilidade e estabelece os limites de latência aceitáveis para as operações mais comuns da plataforma. Tempos de resposta elevados comprometem a experiência de utilização e podem desmotivar a adoção do sistema, particularmente em contextos operacionais urgentes como a marcação de presenças em campo ou a confirmação de uma convocatória.

Os objetivos de desempenho definidos são: as chamadas à API de ações rotineiras devem ser concluídas em no máximo 500ms; as páginas principais da aplicação devem carregar em menos de 3 segundos em redes 4G; e os \textit{dashboards} com dados agregados devem ser renderizados em menos de 5 segundos. Estes valores devem ser verificados em condições de utilização normal e sob carga moderada.

{\small
\rowcolors{2}{rowblue}{white}
\begin{xltabular}{\textwidth}{|l|X|}
\hline
\rowcolor{headerblue}\textcolor{white}{\textbf{Campo}} & \textcolor{white}{\textbf{Descrição}} \\ \hline
\endfirsthead
\multicolumn{2}{c}{{\tablename\ \thetable{} -- Continuação}} \\
\hline
\rowcolor{headerblue}\textcolor{white}{\textbf{Campo}} & \textcolor{white}{\textbf{Descrição}} \\ \hline
\endhead
\hline \multicolumn{2}{|r|}{\textit{Continua na próxima página...}} \\ \hline \endfoot
\hline
\caption{Especificação do Requisito Não Funcional: RNF-06}
\label{tab:rnf-06} \\
\endlastfoot
\textbf{ID} & RNF-06 \\ \hline
\textbf{Nome do Requisito} & Tempos de Resposta Rápidos (Performance) \\ \hline
\textbf{Categoria} & Desempenho e Escalabilidade \\ \hline
\textbf{Descrição} & O sistema deve garantir tempos de carregamento curtos para as ações rotineiras, assegurando uma experiência de utilização fluida mesmo em condições de conectividade móvel, o que é essencial dado o contexto operacional em campo de muitos utilizadores da plataforma. \\ \hline
\textbf{Prioridade} & Alta \\ \hline
\textbf{Origem} & Equipa de Desenvolvimento \\ \hline
\textbf{Analista} & Daniel Parente \\ \hline
\textbf{Critério de Verificação} & Chamadas à API de ações rotineiras concluídas em $\leq$500ms; páginas principais carregadas em $<$3s em rede 4G; \textit{dashboards} renderizados em $<$5s. \\
\end{xltabular}
}
\rowcolors{2}{white}{white}

% ---------------------------------------------------------
% REQUISITO RNF-07
% ---------------------------------------------------------

O requisito RNF-07, denominado Suporte a Picos de Acesso (Escalabilidade), enquadra-se na categoria de Desempenho e Escalabilidade e define a capacidade da arquitetura para absorver aumentos súbitos de tráfego sem degradação do serviço. Eventos como o início de uma época desportiva, a abertura de inscrições ou a divulgação de resultados podem gerar picos de acessos simultâneos que ultrapassam o volume habitual.

A arquitetura deve ser dimensionada para suportar no mínimo 100 utilizadores simultâneos sem perda de desempenho percetível. A base de dados deve ser otimizada com índices adequados para as consultas mais frequentes, e a infraestrutura de alojamento em nuvem deve permitir o escalonamento automático de recursos em resposta a aumentos de carga.

{\small
\rowcolors{2}{rowblue}{white}
\begin{xltabular}{\textwidth}{|l|X|}
\hline
\rowcolor{headerblue}\textcolor{white}{\textbf{Campo}} & \textcolor{white}{\textbf{Descrição}} \\ \hline
\endfirsthead
\multicolumn{2}{c}{{\tablename\ \thetable{} -- Continuação}} \\
\hline
\rowcolor{headerblue}\textcolor{white}{\textbf{Campo}} & \textcolor{white}{\textbf{Descrição}} \\ \hline
\endhead
\hline \multicolumn{2}{|r|}{\textit{Continua na próxima página...}} \\ \hline \endfoot
\hline
\caption{Especificação do Requisito Não Funcional: RNF-07}
\label{tab:rnf-07} \\
\endlastfoot
\textbf{ID} & RNF-07 \\ \hline
\textbf{Nome do Requisito} & Suporte a Picos de Acesso (Escalabilidade) \\ \hline
\textbf{Categoria} & Desempenho e Escalabilidade \\ \hline
\textbf{Descrição} & A arquitetura do sistema deve ser capaz de lidar com aumentos súbitos de tráfego (ex: início de época, abertura de inscrições) sem degradação percetível do serviço, suportando picos de utilização simultânea através de escalonamento automático de recursos. \\ \hline
\textbf{Prioridade} & Média \\ \hline
\textbf{Origem} & Equipa de Desenvolvimento \\ \hline
\textbf{Analista} & Daniel Parente \\ \hline
\textbf{Critério de Verificação} & Sistema suporta $\geq$100 utilizadores simultâneos sem degradação de desempenho; base de dados otimizada com índices nas colunas de filtragem frequente; infraestrutura \textit{cloud} com escalonamento automático configurado. \\
\end{xltabular}
}
\rowcolors{2}{white}{white}

% ---------------------------------------------------------
% REQUISITO RNF-08
% ---------------------------------------------------------

O requisito RNF-08, denominado Alta Disponibilidade (\textit{Uptime}), enquadra-se na categoria de Disponibilidade e Fiabilidade e estabelece os objetivos de continuidade do serviço da plataforma. Para uma organização desportiva que utiliza a plataforma para gerir treinos, convocatórias e pagamentos em ciclo contínuo, paragens não planeadas representam uma disrupção operacional significativa que pode afetar a confiança dos associados.

O sistema deve garantir uma disponibilidade mínima de 99,9\% ao longo do ano civil, o que corresponde a um máximo de aproximadamente 8,7 horas de indisponibilidade acumulada por ano. As janelas de manutenção programada devem ser comunicadas com pelo menos 24 horas de antecedência, e o sistema de monitorização em tempo real deve alertar a equipa técnica de forma imediata em caso de anomalia.

{\small
\rowcolors{2}{rowblue}{white}
\begin{xltabular}{\textwidth}{|l|X|}
\hline
\rowcolor{headerblue}\textcolor{white}{\textbf{Campo}} & \textcolor{white}{\textbf{Descrição}} \\ \hline
\endfirsthead
\multicolumn{2}{c}{{\tablename\ \thetable{} -- Continuação}} \\
\hline
\rowcolor{headerblue}\textcolor{white}{\textbf{Campo}} & \textcolor{white}{\textbf{Descrição}} \\ \hline
\endhead
\hline \multicolumn{2}{|r|}{\textit{Continua na próxima página...}} \\ \hline \endfoot
\hline
\caption{Especificação do Requisito Não Funcional: RNF-08}
\label{tab:rnf-08} \\
\endlastfoot
\textbf{ID} & RNF-08 \\ \hline
\textbf{Nome do Requisito} & Alta Disponibilidade (Uptime) \\ \hline
\textbf{Categoria} & Disponibilidade e Fiabilidade \\ \hline
\textbf{Descrição} & O sistema deve estar disponível de forma contínua, minimizando paragens não planeadas, de modo a garantir que as operações do clube não sejam interrompidas por falhas da plataforma em momentos críticos da atividade desportiva e administrativa. \\ \hline
\textbf{Prioridade} & Alta \\ \hline
\textbf{Origem} & Entrevista --- Gerência do CAP (Reunião 1) \\ \hline
\textbf{Analista} & Tiago Martins \\ \hline
\textbf{Critério de Verificação} & \textit{Uptime} mínimo de 99,9\% anual; manutenção programada comunicada com 24h de antecedência; sistema de monitorização em tempo real com alertas automáticos para falhas. \\
\end{xltabular}
}
\rowcolors{2}{white}{white}

% ---------------------------------------------------------
% REQUISITO RNF-09
% ---------------------------------------------------------

O requisito RNF-09, designado Cópias de Segurança (\textit{Backups}), enquadra-se na categoria de Disponibilidade e Fiabilidade e define a política de proteção dos dados da plataforma contra perdas acidentais, falhas de \textit{hardware} ou incidentes de segurança. Numa plataforma que centraliza dados desportivos, financeiros e de saúde, a perda de informação seria catastrófica para o clube e para os seus associados.

O sistema deve executar cópias de segurança automáticas diárias da base de dados e dos ficheiros armazenados, guardando-as num local geograficamente distinto do servidor principal. Todos os \textit{backups} devem ser encriptados em conformidade com os requisitos legais. Em caso de falha catastrófica, o Objetivo de Tempo de Recuperação (RTO) máximo admissível é de quatro horas.

{\small
\rowcolors{2}{rowblue}{white}
\begin{xltabular}{\textwidth}{|l|X|}
\hline
\rowcolor{headerblue}\textcolor{white}{\textbf{Campo}} & \textcolor{white}{\textbf{Descrição}} \\ \hline
\endfirsthead
\multicolumn{2}{c}{{\tablename\ \thetable{} -- Continuação}} \\
\hline
\rowcolor{headerblue}\textcolor{white}{\textbf{Campo}} & \textcolor{white}{\textbf{Descrição}} \\ \hline
\endhead
\hline \multicolumn{2}{|r|}{\textit{Continua na próxima página...}} \\ \hline \endfoot
\hline
\caption{Especificação do Requisito Não Funcional: RNF-09}
\label{tab:rnf-09} \\
\endlastfoot
\textbf{ID} & RNF-09 \\ \hline
\textbf{Nome do Requisito} & Cópias de Segurança (Backups) \\ \hline
\textbf{Categoria} & Disponibilidade e Fiabilidade \\ \hline
\textbf{Descrição} & O sistema deve realizar cópias de segurança automáticas e periódicas da base de dados e de todos os ficheiros armazenados, guardando-as num local distinto do servidor principal, de forma encriptada e em conformidade com os requisitos legais de proteção de dados. \\ \hline
\textbf{Prioridade} & Alta \\ \hline
\textbf{Origem} & Equipa de Desenvolvimento \\ \hline
\textbf{Analista} & Daniel Parente \\ \hline
\textbf{Critério de Verificação} & \textit{Backups} automáticos diários verificáveis; ficheiros armazenados de forma encriptada em localização geográfica distinta; Objetivo de Tempo de Recuperação (RTO) máximo de 4 horas. \\
\end{xltabular}
}
\rowcolors{2}{white}{white}

% ---------------------------------------------------------
% REQUISITO RNF-10
% ---------------------------------------------------------

O requisito RNF-10, denominado Integração com \textit{Gateways} de Pagamento, enquadra-se na categoria de Integração e Interoperabilidade e define os requisitos técnicos para a ligação da plataforma a processadores de pagamento externos. A aceitação de pagamentos via Multibanco e MBWay é essencial para maximizar a comodidade dos associados na regularização das suas quotas e para reduzir a carga administrativa da tesouraria.

A integração deve ser realizada através das APIs dos operadores certificados (Ifthenpay, SIBS ou Easypay), garantindo que a geração de uma referência de pagamento é concluída em menos de 3 segundos. O sistema deve processar de forma fiável os \textit{callbacks} de confirmação de pagamento enviados pelo \textit{gateway}, atualizando automaticamente o estado da conta corrente do associado sem intervenção manual.

{\small
\rowcolors{2}{rowblue}{white}
\begin{xltabular}{\textwidth}{|l|X|}
\hline
\rowcolor{headerblue}\textcolor{white}{\textbf{Campo}} & \textcolor{white}{\textbf{Descrição}} \\ \hline
\endfirsthead
\multicolumn{2}{c}{{\tablename\ \thetable{} -- Continuação}} \\
\hline
\rowcolor{headerblue}\textcolor{white}{\textbf{Campo}} & \textcolor{white}{\textbf{Descrição}} \\ \hline
\endhead
\hline \multicolumn{2}{|r|}{\textit{Continua na próxima página...}} \\ \hline \endfoot
\hline
\caption{Especificação do Requisito Não Funcional: RNF-10}
\label{tab:rnf-10} \\
\endlastfoot
\textbf{ID} & RNF-10 \\ \hline
\textbf{Nome do Requisito} & Integração com Gateways de Pagamento \\ \hline
\textbf{Categoria} & Integração e Interoperabilidade \\ \hline
\textbf{Descrição} & O sistema deve estabelecer uma ligação robusta via API com processadores de pagamento externos (Ifthenpay, SIBS, Easypay) para a geração de referências Multibanco/MBWay e para o processamento fiável de \textit{callbacks} de confirmação de pagamento. \\ \hline
\textbf{Prioridade} & Alta \\ \hline
\textbf{Origem} & Análise Documental (Fatura/Recibo de Quota) \\ \hline
\textbf{Analista} & Pedro Ferreira \\ \hline
\textbf{Critério de Verificação} & Referência de pagamento gerada em $\leq$3 segundos; \textit{callbacks} do \textit{gateway} processados e conta corrente atualizada automaticamente; testes de integração com ambiente de \textit{sandbox} do operador. \\
\end{xltabular}
}
\rowcolors{2}{white}{white}

% ---------------------------------------------------------
% REQUISITO RNF-11
% ---------------------------------------------------------

O requisito RNF-11, denominado Exportação de Dados e Normalização, enquadra-se na categoria de Integração e Interoperabilidade e estabelece as normas para a geração de ficheiros em formatos universais que permitam a interoperabilidade com sistemas externos. A plataforma deve ser capaz de produzir ficheiros fiscais (SAF-T), relatórios financeiros (Excel/CSV) e documentos formatados (PDF) que possam ser utilizados por contabilistas, federações e outros intervenientes externos.

Os ficheiros SAF-T gerados devem passar na validação do validador oficial da Autoridade Tributária. Os ficheiros Excel e CSV devem ser codificados em UTF-8 para garantir a correta representação de caracteres portugueses. Os documentos PDF devem incluir o logótipo do clube, paginação e formatação adequada para impressão ou arquivo digital.

{\small
\rowcolors{2}{rowblue}{white}
\begin{xltabular}{\textwidth}{|l|X|}
\hline
\rowcolor{headerblue}\textcolor{white}{\textbf{Campo}} & \textcolor{white}{\textbf{Descrição}} \\ \hline
\endfirsthead
\multicolumn{2}{c}{{\tablename\ \thetable{} -- Continuação}} \\
\hline
\rowcolor{headerblue}\textcolor{white}{\textbf{Campo}} & \textcolor{white}{\textbf{Descrição}} \\ \hline
\endhead
\hline \multicolumn{2}{|r|}{\textit{Continua na próxima página...}} \\ \hline \endfoot
\hline
\caption{Especificação do Requisito Não Funcional: RNF-11}
\label{tab:rnf-11} \\
\endlastfoot
\textbf{ID} & RNF-11 \\ \hline
\textbf{Nome do Requisito} & Exportação de Dados e Normalização \\ \hline
\textbf{Categoria} & Integração e Interoperabilidade \\ \hline
\textbf{Descrição} & O sistema deve suportar a geração e exportação de ficheiros em formatos universais (SAF-T XML, Excel, CSV, PDF), garantindo a interoperabilidade com sistemas externos utilizados por contabilistas, federações desportivas e outros intervenientes da organização. \\ \hline
\textbf{Prioridade} & Alta \\ \hline
\textbf{Origem} & Entrevista --- Gerência do CAP (Reunião 1) \\ \hline
\textbf{Analista} & Tiago Martins \\ \hline
\textbf{Critério de Verificação} & Ficheiro SAF-T passa na validação do validador oficial da AT; Excel/CSV exportados em codificação UTF-8; PDF gerado com logótipo do clube, paginação e formatação adequada para impressão. \\
\end{xltabular}
}
\rowcolors{2}{white}{white}

% ---------------------------------------------------------
% REQUISITO RNF-12
% ---------------------------------------------------------

O requisito RNF-12, denominado Arquitetura Modular e Escalável, enquadra-se na categoria de Manutenibilidade e Auditoria e estabelece as diretrizes de desenho de \textit{software} que a plataforma deve seguir para garantir a sua sustentabilidade técnica a longo prazo. Um sistema que não respeite princípios de modularidade e separação de responsabilidades torna-se rapidamente difícil de manter, evoluir e adaptar às necessidades em constante mudança de uma organização desportiva.

O código deve ser estruturado segundo os princípios da \textit{Clean Architecture} ou do padrão MVC, com separação clara entre as camadas de apresentação, lógica de negócio e acesso a dados. Esta organização deve permitir a adição de novas funcionalidades sem comprometer a estabilidade do sistema existente, minimizando o risco de regressões e facilitando os processos de revisão de código.

{\small
\rowcolors{2}{rowblue}{white}
\begin{xltabular}{\textwidth}{|l|X|}
\hline
\rowcolor{headerblue}\textcolor{white}{\textbf{Campo}} & \textcolor{white}{\textbf{Descrição}} \\ \hline
\endfirsthead
\multicolumn{2}{c}{{\tablename\ \thetable{} -- Continuação}} \\
\hline
\rowcolor{headerblue}\textcolor{white}{\textbf{Campo}} & \textcolor{white}{\textbf{Descrição}} \\ \hline
\endhead
\hline \multicolumn{2}{|r|}{\textit{Continua na próxima página...}} \\ \hline \endfoot
\hline
\caption{Especificação do Requisito Não Funcional: RNF-12}
\label{tab:rnf-12} \\
\endlastfoot
\textbf{ID} & RNF-12 \\ \hline
\textbf{Nome do Requisito} & Arquitetura Modular e Escalável \\ \hline
\textbf{Categoria} & Manutenibilidade e Auditoria \\ \hline
\textbf{Descrição} & O código da plataforma deve ser desenvolvido com \textit{Clean Architecture} ou padrão MVC, com separação clara de responsabilidades em camadas lógicas, de modo a permitir a adição de novas funcionalidades (modalidades, escalões, integrações) sem comprometer a estabilidade do sistema existente. \\ \hline
\textbf{Prioridade} & Alta \\ \hline
\textbf{Origem} & Equipa de Desenvolvimento \\ \hline
\textbf{Analista} & Daniel Parente \\ \hline
\textbf{Critério de Verificação} & Código documentado e organizado em camadas lógicas distintas; adição de novo escalão ou modalidade realizável sem modificar lógica de negócio existente; ausência de dependências circulares entre módulos. \\
\end{xltabular}
}
\rowcolors{2}{white}{white}

% ---------------------------------------------------------
% REQUISITO RNF-13
% ---------------------------------------------------------

O requisito RNF-13, denominado Registos de Auditoria (\textit{Logs}), enquadra-se na categoria de Manutenibilidade e Auditoria e estabelece a obrigação de manter um registo persistente e imutável de todas as ações críticas realizadas na plataforma. A existência de \textit{logs} de auditoria é fundamental para a rastreabilidade de operações sensíveis, para a resolução de disputas entre utilizadores e para a demonstração de conformidade com os requisitos legais e regulatórios.

Cada entrada no registo de auditoria deve conter, no mínimo, o identificador do utilizador, o \textit{timestamp} com fuso horário, o endereço IP da origem, a descrição da ação realizada e o identificador do recurso afetado. Os \textit{logs} devem ser armazenados de forma que impeça a sua edição ou eliminação por qualquer utilizador, incluindo administradores do sistema, garantindo a sua integridade enquanto registo de evidências.

{\small
\rowcolors{2}{rowblue}{white}
\begin{xltabular}{\textwidth}{|l|X|}
\hline
\rowcolor{headerblue}\textcolor{white}{\textbf{Campo}} & \textcolor{white}{\textbf{Descrição}} \\ \hline
\endfirsthead
\multicolumn{2}{c}{{\tablename\ \thetable{} -- Continuação}} \\
\hline
\rowcolor{headerblue}\textcolor{white}{\textbf{Campo}} & \textcolor{white}{\textbf{Descrição}} \\ \hline
\endhead
\hline \multicolumn{2}{|r|}{\textit{Continua na próxima página...}} \\ \hline \endfoot
\hline
\caption{Especificação do Requisito Não Funcional: RNF-13}
\label{tab:rnf-13} \\
\endlastfoot
\textbf{ID} & RNF-13 \\ \hline
\textbf{Nome do Requisito} & Registos de Auditoria (Logs) \\ \hline
\textbf{Categoria} & Manutenibilidade e Auditoria \\ \hline
\textbf{Descrição} & O sistema deve manter um registo persistente de todas as ações críticas (autenticações, operações financeiras, eliminação de dados, submissão de documentos, alterações de privilégios), garantindo a rastreabilidade completa das operações para fins de auditoria interna e conformidade legal. \\ \hline
\textbf{Prioridade} & Alta \\ \hline
\textbf{Origem} & Equipa de Desenvolvimento \\ \hline
\textbf{Analista} & Daniel Parente \\ \hline
\textbf{Critério de Verificação} & Cada \textit{log} contém: ID do utilizador, \textit{timestamp} com fuso horário, IP de origem, descrição da ação e ID do recurso afetado; \textit{logs} não editáveis nem elimináveis por nenhum perfil de utilizador; retenção mínima de 5 anos. \\
\end{xltabular}
}
\rowcolors{2}{white}{white}




\rowcolors{2}{rowblue}{white}
\begin{table}[H]
\centering
\begin{tabularx}{\textwidth}{|l|X|}
\hline
\rowcolor{headerblue}\textcolor{white}{\textbf{Campo}} & \textcolor{white}{\textbf{Descrição}} \\ \hline
\textbf{ID} & RF-4.1 \\ \hline
\textbf{Nome} & Gestão de Inscrições e Renovações \\ \hline
\textbf{Atores} & Atleta, Encarregado de Educação, Secretaria \\ \hline
\textbf{Descrição} & Processar os pedidos de inscrição e renovação anual dos atletas, garantindo a regularização documental e desportiva. \\ \hline
\textbf{Pré-condições} & Perfil de utilizador criado; época desportiva aberta. \\ \hline
\textbf{Pós-condições} & Atleta inscrito na modalidade/escalão; plano de pagamento de quotas gerado automaticamente. \\ \hline
\textbf{Fluxo Principal} & 
1. O utilizador submete o formulário de inscrição e anexa o documento de identificação e o EMD. \newline
2. O sistema coloca o pedido no estado ``Em Análise''. \newline
3. A secretaria valida a autenticidade e validade dos documentos. \newline
4. A secretaria aprova a inscrição. \newline
5. O sistema ativa o atleta para a época e emite as faturas correspondentes. \\ \hline
\textbf{Fluxo de Exceção} & 
[Documentos Inválidos] (passo 3) \newline
3.1. A secretaria deteta falhas (ex: EMD caducado). \newline
3.2. A secretaria rejeita o pedido e solicita correções via notificação automática. \\ \hline
\end{tabularx}
\caption{Especificação do Requisito Funcional: RF-4.1}
\label{tab:rf-4.1}
\end{table}
\rowcolors{2}{white}{white}

A salvaguarda da saúde e o acompanhamento do desenvolvimento físico são pilares do dever de zelo do clube. Os requisitos seguintes detalham o controlo de indisponibilidades médicas e a monitorização biométrica.

\rowcolors{2}{rowblue}{white}
\begin{table}[H]
\centering
\begin{tabularx}{\textwidth}{|l|X|}
\hline
\rowcolor{headerblue}\textcolor{white}{\textbf{Campo}} & \textcolor{white}{\textbf{Descrição}} \\ \hline
\textbf{ID} & RF-4.2 \\ \hline
\textbf{Nome} & Gestão de Lesões e Recuperação \\ \hline
\textbf{Atores} & Atleta, Treinador, Departamento Médico \\ \hline
\textbf{Descrição} & Registo de queixas e lesões, com bloqueio automático de convocação e acompanhamento clínico. \\ \hline
\textbf{Métrica de Aceitação} & Impedimento de convocação imediato após registo de lesão; histórico clínico visível apenas para perfis autorizados (RGPD). \\ \hline
\end{tabularx}
\caption{Especificação do Requisito Funcional: RF-4.2}
\label{tab:rf-4.2}
\end{table}
\rowcolors{2}{white}{white}

\rowcolors{2}{rowblue}{white}
\begin{table}[H]
\centering
\begin{tabularx}{\textwidth}{|l|X|}
\hline
\rowcolor{headerblue}\textcolor{white}{\textbf{Campo}} & \textcolor{white}{\textbf{Descrição}} \\ \hline
\textbf{ID} & RF-4.3 \\ \hline
\textbf{Nome} & Acompanhamento Físico e Biometria \\ \hline
\textbf{Atores} & Treinador, Atleta \\ \hline
\textbf{Descrição} & Registo periódico de peso, altura e cálculo automático de IMC para análise de crescimento. \\ \hline
\textbf{Métrica de Aceitação} & Gráficos de evolução gerados automaticamente no perfil do atleta. \\ \hline
\end{tabularx}
\caption{Especificação do Requisito Funcional: RF-4.3}
\label{tab:rf-4.3}
\end{table}
\rowcolors{2}{white}{white}

\paragraph{Subsistema de Análise e Desempenho}

A competitividade e a melhoria técnica são suportadas por dados objetivos de performance, permitindo a comparação entre pares e a definição de metas.

\rowcolors{2}{rowblue}{white}
\begin{table}[H]
\centering
\begin{tabularx}{\textwidth}{|l|X|}
\hline
\rowcolor{headerblue}\textcolor{white}{\textbf{Campo}} & \textcolor{white}{\textbf{Descrição}} \\ \hline
\textbf{ID} & RF-5.1 \\ \hline
\textbf{Nome} & Painéis de Desempenho (Dashboards) \\ \hline
\textbf{Atores} & Atleta, Treinador, Gerência \\ \hline
\textbf{Descrição} & Visualização gráfica de estatísticas de jogo, assiduidade e evolução física. \\ \hline
\textbf{Métrica de Aceitação} & Disponibilização de gráficos radar para comparação de atletas da mesma posição. \\ \hline
\end{tabularx}
\caption{Especificação do Requisito Funcional: RF-5.1}
\label{tab:rf-5.1}
\end{table}
\rowcolors{2}{white}{white}

\rowcolors{2}{rowblue}{white}
\begin{table}[H]
\centering
\begin{tabularx}{\textwidth}{|l|X|}
\hline
\rowcolor{headerblue}\textcolor{white}{\textbf{Campo}} & \textcolor{white}{\textbf{Descrição}} \\ \hline
\textbf{ID} & RF-5.2 \\ \hline
\textbf{Nome} & Avaliação Qualitativa e Objetivos \\ \hline
\textbf{Atores} & Treinador, Coordenador \\ \hline
\textbf{Descrição} & Atribuição de notas de empenho e técnica no final dos treinos e definição de OKRs individuais. \\ \hline
\textbf{Métrica de Aceitação} & Interface de feedback rápido pós-treino; relatório mensal de objetivos enviado ao atleta. \\ \hline
\end{tabularx}
\caption{Especificação do Requisito Funcional: RF-5.2}
\label{tab:rf-5.2}
\end{table}
\rowcolors{2}{white}{white}

O requisito que se segue, designado por Gestão de Inscrições e Renovações (RF-4.1), dá continuidade ao processo de estruturação administrativa do clube. Uma vez que os perfis estão criados e os atletas alocados aos respetivos plantéis, é estritamente necessário garantir que a situação legal e desportiva de cada praticante está regularizada para a época em curso. Este requisito determina que os atletas, ou os seus encarregados de educação, podem submeter o pedido de inscrição ou renovação através da plataforma, anexando a documentação desportiva e clínica exigida.

Tratando-se de um processo que envolve a submissão e verificação de ficheiros confidenciais, foi estipulado que a documentação submetida transita para um estado de análise e que a validação e aprovação final destas inscrições apenas pode ser levada a cabo pelos funcionários da secretaria do clube. A conclusão e o cumprimento deste requisito revelam-se de extrema importância para o sistema, uma vez que apenas os atletas com a inscrição devidamente aprovada poderão ser oficialmente inseridos nas convocatórias e participar em competições, servindo também de ponto de partida para a geração automática do plano de pagamento de quotas associado a essa época.

\rowcolors{2}{rowblue}{white}
\begin{table}[H]
\centering
\begin{tabularx}{\textwidth}{|l|X|}
\hline
\rowcolor{headerblue}\textcolor{white}{\textbf{Campo}} & \textcolor{white}{\textbf{Descrição}} \\ \hline
\textbf{ID} & RF-2.2 \\ \hline
\textbf{Nome} & Notificações e Alertas \\ \hline
\textbf{Atores} & Sistema, Atleta, Treinador, Encarregado \\ \hline
\textbf{Descrição} & Garantir a comunicação proativa de eventos críticos através de avisos automáticos multicanal. \\ \hline
\textbf{Pré-condições} & Utilizador tem dados de contacto (email/telemóvel) no perfil; evento despoletado. \\ \hline
\textbf{Pós-condições} & Informação entregue ao utilizador; registo de log de envio. \\ \hline
\textbf{Fluxo Principal} & 
1. O sistema deteta um gatilho de negócio (ex: alteração de horário de treino). \newline
2. O sistema recupera a lista de utilizadores afetados e as suas preferências de receção. \newline
3. O sistema formata a mensagem utilizando os modelos de notificação. \newline
4. O sistema processa o envio imediato via correio eletrónico ou notificação \textit{push}. \\ \hline
\textbf{Fluxo de Exceção} & 
[Falha no serviço de envio] (passo 4) \newline
4.1. O gateway de terceiros (ex: SendGrid) devolve um erro. \newline
4.2. O sistema enfileira a notificação para nova tentativa posterior. \\ \hline
\end{tabularx}
\caption{Especificação do Requisito Funcional: RF-2.2}
\label{tab:rf-2.2}
\end{table}
\rowcolors{2}{white}{white}

O quarto requisito funcional, designado por Notificações e Alertas (RF-2.2), assume um papel transversal e de suporte a toda a dinâmica da plataforma. Este requisito surge da necessidade de garantir que o sistema gera e envia avisos automáticos aos diversos utilizadores (como atletas, encarregados de educação, treinadores e staff) sempre que ocorram eventos que exijam a sua atenção. Entre estes eventos incluem-se a publicação de novas convocatórias, alterações de última hora em horários ou locais de treino, e avisos de quotas ou pagamentos em atraso.

Para que as notificações sejam emitidas com sucesso, verifica-se como pré-condição que o utilizador destinatário tenha os seus dados de contacto devidamente preenchidos no perfil e que o evento que despoleta o aviso tenha sido acionado no sistema. Assim que a ação ocorre, o sistema atua de forma proativa para cumprir a pós-condição: a informação chega imediatamente ao utilizador, seja através de alertas visuais internos na própria interface da aplicação ou através de comunicações externas, como o correio eletrónico. A relevância deste requisito para a aplicação é bastante elevada, dado que assegura uma comunicação fluida, contínua e atempada entre a estrutura do clube e os seus membros, mitigando falhas de informação no dia a dia.

\rowcolors{2}{rowblue}{white}
\begin{table}[H]
\centering
\begin{tabularx}{\textwidth}{|l|X|}
\hline
\rowcolor{headerblue}\textcolor{white}{\textbf{Campo}} & \textcolor{white}{\textbf{Descrição}} \\ \hline
\textbf{ID} & RF-3.4 \\ \hline
\textbf{Nome} & Gestão de Convocatórias \\ \hline
\textbf{Atores} & Treinador, Atleta, Encarregado de Educação \\ \hline
\textbf{Descrição} & Definição e comunicação da equipa selecionada para eventos competitivos ou amigáveis. \\ \hline
\textbf{Pré-condições} & Treinador autenticado; evento criado no calendário. \\ \hline
\textbf{Pós-condições} & Lista de convocados persistida; atletas notificados; estado da convocatória alterado para ``Publicada''. \\ \hline
\textbf{Fluxo Principal} & 
1. O treinador seleciona a equipa e o jogo alvo. \newline
2. O treinador escolhe os atletas da lista de disponíveis (aptos). \newline
3. O treinador submete a seleção para publicação oficial. \newline
4. O sistema gera os alertas de convocatória para os envolvidos. \newline
5. O sistema recolhe as confirmações de presença em tempo real. \\ \hline
\textbf{Fluxo de Exceção} & 
[Atleta Inapto] (passo 2) \newline
2.1. O treinador tenta convocar um atleta com EMD caducado ou quotas suspensas. \newline
2.2. O sistema bloqueia a seleção e indica o motivo da inaptidão. \\ \hline
\end{tabularx}
\caption{Especificação do Requisito Funcional: RF-3.4}
\label{tab:rf-3.4}
\end{table}
\rowcolors{2}{white}{white}

A gestão desportiva é complementada pelo planeamento operacional e registo de performance quotidiana. Os requisitos seguintes detalham o controlo de horários, assiduidade, planos de treino e o registo disciplinar e de jogo.

\rowcolors{2}{rowblue}{white}
\begin{table}[H]
\centering
\begin{tabularx}{\textwidth}{|l|X|}
\hline
\rowcolor{headerblue}\textcolor{white}{\textbf{Campo}} & \textcolor{white}{\textbf{Descrição}} \\ \hline
\textbf{ID} & RF-3.2 \\ \hline
\textbf{Nome} & Planeamento de Horários e Treinos \\ \hline
\textbf{Atores} & Treinador, Gerência, Atleta \\ \hline
\textbf{Descrição} & Permitir a definição e alteração dinâmica do calendário de treinos, com atualização imediata para os atletas. \\ \hline
\textbf{Métrica de Aceitação} & Atualização do calendário em tempo real (< 1s) após gravação; envio automático de notificações de alteração. \\ \hline
\end{tabularx}
\caption{Especificação do Requisito Funcional: RF-3.2}
\label{tab:rf-3.2}
\end{table}
\rowcolors{2}{white}{white}

\rowcolors{2}{rowblue}{white}
\begin{table}[H]
\centering
\begin{tabularx}{\textwidth}{|l|X|}
\hline
\rowcolor{headerblue}\textcolor{white}{\textbf{Campo}} & \textcolor{white}{\textbf{Descrição}} \\ \hline
\textbf{ID} & RF-3.3 \\ \hline
\textbf{Nome} & Gestão de Presenças e Faltas \\ \hline
\textbf{Atores} & Treinador, Atleta \\ \hline
\textbf{Descrição} & Registo sistemático da assiduidade dos atletas em cada sessão de treino para fins de avaliação e seguro. \\ \hline
\textbf{Métrica de Aceitação} & Interface de chamada rápida disponível na App do Treinador; histórico de presenças consultável no perfil do atleta. \\ \hline
\end{tabularx}
\caption{Especificação do Requisito Funcional: RF-3.3}
\label{tab:rf-3.3}
\end{table}
\rowcolors{2}{white}{white}

\rowcolors{2}{rowblue}{white}
\begin{table}[H]
\centering
\begin{tabularx}{\textwidth}{|l|X|}
\hline
\rowcolor{headerblue}\textcolor{white}{\textbf{Campo}} & \textcolor{white}{\textbf{Descrição}} \\ \hline
\textbf{ID} & RF-3.5 \\ \hline
\textbf{Nome} & Planos de Treino e Recursos Táticos \\ \hline
\textbf{Atores} & Treinador, Atleta \\ \hline
\textbf{Descrição} & Disponibilização de conteúdos técnicos (vídeos, diagramas) e planos de exercício individualizados. \\ \hline
\textbf{Métrica de Aceitação} & Suporte para anexação de ficheiros multimédia; feed de recursos filtrado por equipa e posição. \\ \hline
\end{tabularx}
\caption{Especificação do Requisito Funcional: RF-3.5}
\label{tab:rf-3.5}
\end{table}
\rowcolors{2}{white}{white}

\rowcolors{2}{rowblue}{white}
\begin{table}[H]
\centering
\begin{tabularx}{\textwidth}{|l|X|}
\hline
\rowcolor{headerblue}\textcolor{white}{\textbf{Campo}} & \textcolor{white}{\textbf{Descrição}} \\ \hline
\textbf{ID} & RF-3.6 \\ \hline
\textbf{Nome} & Registo de Ocorrências e Disciplina \\ \hline
\textbf{Atores} & Treinador, Secretaria, Gerência \\ \hline
\textbf{Descrição} & Registo formal de eventos de jogo (golos, cartões) e incidências disciplinares (atrasos, suspensões). \\ \hline
\textbf{Métrica de Aceitação} & Consolidação automática das ocorrências no palmarés do atleta e da equipa; alertas de suspensão automáticos. \\ \hline
\end{tabularx}
\caption{Especificação do Requisito Funcional: RF-3.6}
\label{tab:rf-3.6}
\end{table}
\rowcolors{2}{white}{white}

O sexto requisito funcional, designado por Gestão de Convocatórias (RF-3.4), dá continuidade à vertente de Gestão Desportiva da plataforma. Este requisito surge da necessidade da equipa técnica definir e comunicar oficialmente quais os atletas selecionados para representar o clube num determinado jogo, seja ele de caráter oficial ou amigável.

Para a realização desta operação, verifica-se como pré-condição que o utilizador se encontre autenticado com o perfil de Treinador ou Coordenador Técnico e que o evento correspondente ao jogo já esteja criado no calendário do sistema. Durante o processo de seleção, a interface deve auxiliar a decisão do treinador, cruzando os dados do plantel com as informações de assiduidade aos treinos e o histórico clínico (para evitar a convocação de atletas lesionados ou a cumprir castigo). Uma vez confirmada e submetida a seleção, cumpre-se a pós-condição: a lista de convocados é guardada na base de dados de forma persistente e o sistema aciona de imediato o envio de alertas automáticos aos atletas selecionados e aos respetivos encarregados de educação.

A relevância deste requisito para a aplicação é bastante elevada, uma vez que representa o culminar do ciclo de trabalho semanal da equipa. A desmaterialização das convocatórias garante que a informação chega de forma rápida e inequívoca aos intervenientes envolvidos, minimizando falhas de comunicação e facilitando toda a organização logística para os dias de jogo.


\rowcolors{2}{rowblue}{white}
\begin{table}[H]
\centering
\begin{tabularx}{\textwidth}{|l|X|}
\hline
\rowcolor{headerblue}\textcolor{white}{\textbf{Campo}} & \textcolor{white}{\textbf{Descrição}} \\ \hline
\textbf{ID} & RF-7.1 \\ \hline
\textbf{Nome} & Gestão de Reserva de Espaços \\ \hline
\textbf{Atores} & Secretaria, Coordenadores, Treinadores \\ \hline
\textbf{Descrição} & Agendamento e controlo de ocupação das infraestruturas físicas do clube. \\ \hline
\textbf{Pré-condições} & Utilizador autenticado com permissões logísticas; instalações mapeadas. \\ \hline
\textbf{Pós-condições} & Reserva registada e visível no calendário global; prevenção de duplicidade. \\ \hline
\textbf{Fluxo Principal} & 
1. O utilizador consulta a disponibilidade no calendário interativo. \newline
2. O utilizador seleciona o espaço e define a janela temporal. \newline
3. O sistema valida se o recurso está livre para o período solicitado. \newline
4. O sistema confirma a reserva e atualiza o mapa público de instalações. \\ \hline
\textbf{Fluxo de Exceção} & 
[Conflito de Horário] (passo 3) \newline
3.1. O sistema deteta que o espaço já está alocado a outra equipa ou em manutenção. \newline
3.2. O sistema emite um aviso e apresenta sugestões de horários adjacentes livres. \\ \hline
\end{tabularx}
\caption{Especificação do Requisito Funcional: RF-7.3}
\label{tab:rf-7.3}
\end{table}
\rowcolors{2}{white}{white}

A gestão financeira e logística constitui a espinha dorsal da sustentabilidade do CAP. O motor de faturação integrado (RF-6.1) elimina os erros de processamento manual identificados na fase de diagnóstico, garantindo que o clube mantém um fluxo de caixa estável e auditável. A exportação automática do SAF-T (RF-6.3) assegura a conformidade legal sem custos administrativos adicionais. Na vertente logística, a centralização das reservas (RF-7.1) e o controlo de stock (RF-7.3) permitem uma otimização sem precedentes dos ativos físicos do clube, desde o relvado até ao fardamento oficial.

\rowcolors{2}{rowblue}{white}
\begin{table}[H]
\centering
\begin{tabularx}{\textwidth}{|l|X|}
\hline
\rowcolor{headerblue}\textcolor{white}{\textbf{Campo}} & \textcolor{white}{\textbf{Descrição}} \\ \hline
\textbf{ID} & RF-7.1 \\ \hline
\textbf{Nome} & Gestão de Reserva de Espaços \\ \hline
\textbf{Atores} & Secretaria, Coordenadores, Treinadores \\ \hline
\textbf{Descrição} & Agendamento e controlo de ocupação das infraestruturas físicas do clube. \\ \hline
\textbf{Pré-condições} & Utilizador autenticado com permissões logísticas; instalações mapeadas. \\ \hline
\textbf{Pós-condições} & Reserva registada e visível no calendário global; prevenção de duplicidade. \\ \hline
\textbf{Fluxo Principal} & 
1. O utilizador consulta a disponibilidade no calendário interativo. \newline
2. O utilizador seleciona o espaço e define a janela temporal. \newline
3. O sistema valida se o recurso está livre para o período solicitado. \newline
4. O sistema confirma a reserva e atualiza o mapa público de instalações. \\ \hline
\textbf{Fluxo de Exceção} & 
[Conflito de Horário] (passo 3) \newline
3.1. O sistema deteta que o espaço já está alocado a outra equipa ou em manutenção. \newline
3.2. O sistema emite um aviso e apresenta sugestões de horários adjacentes livres. \\ \hline
\end{tabularx}
\caption{Especificação do Requisito Funcional: RF-7.1}
\label{tab:rf-7.1}
\end{table}
\rowcolors{2}{white}{white}

O décimo terceiro requisito funcional, designado por Gestão de Reserva de Espaços e Instalações Desportivas (RF-7.1), insere-se na área de Gestão Logística e Administrativa da plataforma. Este requisito surge da necessidade da estrutura do clube, nomeadamente a secretaria e os coordenadores técnicos, gerirem de forma eficiente a ocupação dos campos, pavilhões, balneários e ginásios, evitando sobreposições de horários entre os diferentes escalões e equipas.

Para que a alocação ou reserva de uma instalação seja efetuada, verifica-se como pré-condição que o utilizador responsável se encontre autenticado no sistema com o perfil de Coordenador Técnico, Treinador ou Administrador, e que os espaços físicos disponíveis já estejam previamente mapeados e configurados na base de dados.

Durante a operação, o utilizador acede ao calendário interativo de ocupação, seleciona o espaço pretendido, define a janela temporal e associa a reserva a uma equipa ou evento específico (como um treino, jogo oficial ou manutenção). O sistema efetua imediatamente uma validação de disponibilidade. Caso detete um conflito de horários, a operação é bloqueada e o utilizador é notificado; em caso de sucesso, cumpre-se a pós-condição: a reserva é registada de forma persistente na base de dados e o calendário global do clube é atualizado, ficando visível para todos os perfis autorizados.

A relevância deste requisito para a aplicação é elevada, uma vez que a otimização dos espaços físicos é fundamental para o funcionamento diário da instituição. A centralização e digitalização deste processo eliminam as redundâncias e os conflitos logísticos, garantindo que o planeamento desportivo decorre sem interrupções.

